\documentclass[11pt, letterpaper]{article}
\usepackage{amsmath, amssymb}
\usepackage{booktabs}
\usepackage{multirow}
\usepackage{geometry}
\usepackage{hyperref}
\usepackage{xcolor}
\usepackage{graphicx}
\usepackage{float}
\usepackage{enumitem}
\usepackage{tcolorbox}
\usepackage{algorithm}
\usepackage{algpseudocode}
\geometry{margin=1in}
\tcbuselibrary{skins, breakable}
\newtcolorbox{agentbox}[1]{title={#1}, colback=blue!3!white, colframe=blue!55!black,
  fonttitle=\bfseries, breakable}
\newtcolorbox{keybox}[1]{title={#1}, colback=green!3!white, colframe=green!45!black,
  fonttitle=\bfseries, breakable}
\newtcolorbox{warnbox}[1]{title={#1}, colback=orange!5!white, colframe=orange!70!black,
  fonttitle=\bfseries, breakable}
\hypersetup{colorlinks=true, linkcolor=blue!50!black}

\title{\textbf{Multi-Agent Receding-Horizon Games for Autonomous Market Participation}\\[4pt]
\large Detailed Case-Study Report --- Distributed Day-Ahead ADMM \& Explicit\\
Real-Time FACET Generalized Nash Equilibrium\\
\normalsize Green-H$_2$ coalition, ERCOT 2025, $H=4$ (1-hour) receding horizon, PCC band $[100,900]$ kW}
\author{Parth Brahmbhatt --- University of Wisconsin--Madison}
\date{\today}

\begin{document}
\maketitle
\tableofcontents
\bigskip\hrule\bigskip

\begin{abstract}
Six independently-owned green-hydrogen electrolyzers behind one point-of-common-coupling (PCC) bid
collectively into the ERCOT two-settlement market. \textbf{Both} market stages are solved fully
distributed, over the \emph{same} physical PCC import band $[L_{\min},L_{\max}]=[100,900]$~kW, and
\emph{no agent ever shares its cost model}. The day-ahead (DAM) anchor is set by a distributed
\textbf{ADMM} in which each agent solves only its own 24-hour QP and the sole shared quantity is
the per-hour coalition import. Real-time dispatch (RTM, 15-min) is a \textbf{receding-horizon}
($H=4$) parametric generalized Nash equilibrium (GNE) solved online by \textbf{FACET}: each agent
locates its own precomputed critical region by warm-started point-location over its facet-neighbour
graph (never enumerating the $K^N\!\approx\!1.1\times10^{21}$ combinations), the coalition combination
is made self-consistent, and the \emph{variational} GNE (= social optimum) is recovered by a
potential-minimizing solve that remains valid when the coupling binds. All six agents are distinct
in converter rating and efficiency, so each carries its own offline map. Over two independent ERCOT
weeks (1--7 April, 7--13 July 2025; $-\$1$ to $\$3553$/MWh) FACET clears \textbf{99.3\%} of steps
\emph{iteration-free} --- a \textbf{median of one inter-agent round per step} against the $67$ a
tuned distributed ADMM needs for the same equilibrium ($1.8\%$ of its communication over the two
weeks), at a median \textbf{3.9~ms} against a $900$~s dispatch interval --- meeting the weekly
H$_2$ target at \textbf{140--141\%}, holding renewable curtailment to \textbf{5.8--6.8\%} of
available generation, and matching the centralized equilibrium to $\sim\!10^{-4}$~kW in the
interior. Against ADMM head-to-head, identically initialised and judged by the same metric, FACET
is $\sim\!34\times$ faster at the median step; the robust, implementation-independent claim, though,
is the communication ($1$ vs $67$ rounds). On the $0.7\%$ of steps at the degenerate coupling
ceiling FACET returns a certified but not-necessarily-\emph{variational} GNE (bounded to
$\sim\!2\%$ of the PCC band), and a warm-started distributed ADMM fallback ($\sim\!27$ rounds) is
the only iteration. This report documents each agent's optimization problem, the distributed DAM
and the online FACET clearing step by step, and all measured results including the solver
comparison, solve times and curtailment. \S\ref{sec:limit} reports honestly that the residual
degeneracy --- not the fleet composition --- is what bounds the iteration-free rate.
\end{abstract}

\section{Nomenclature}\label{sec:nomen}
% ─────────────────────────────────────────────────────────────────────────────
Every symbol and every term of art used in this report is defined here, and used nowhere in a looser
sense.

\begin{keybox}{Terms of art}
\begin{description}[leftmargin=0pt,labelindent=0pt,style=nextline,nosep,itemsep=2pt]
\item[Generalized Nash equilibrium (GNE)] A collective decision at which no agent can lower its own
cost by unilaterally changing its own decision. ``Generalized'' means the agents' \emph{feasible sets}
are coupled --- here by the PCC band --- not merely their costs.
\item[Variational GNE (v-GNE)] The particular GNE at which the Lagrange multipliers of the coupling
constraint are equal across agents. For a potential game it coincides with the minimizer of the
potential, i.e.\ the \emph{social optimum}. When several GNE exist, this is the one we select.
\item[Potential game] A game admitting a single function $\Phi$ (the \emph{potential}) whose gradient
reproduces every agent's cost gradient. Consequence used throughout: the v-GNE is
$\arg\min\Phi$ over the joint feasible set.
\item[Multiparametric QP (mpQP)] A quadratic program solved \emph{once, offline}, symbolically in its
parameters, yielding the optimizer as an explicit piecewise-affine function of those parameters.
\item[Critical region (CR)] One piece of that piecewise-affine function: a polytope in parameter space
inside which the optimizer is a single affine function of the parameters. Written $E_it_i\le f_i$
with law $x_i=G_it_i+g_i$.
\item[Combination] A choice of one CR per agent, $\mathcal C=(r_0,\dots,r_5)$. There are
$K^N\approx1.1\times10^{21}$ of them; we never enumerate them.
\item[Point location] Determining which CR contains a given parameter point, by the containment test
$E_it_i\le f_i$. Pure arithmetic: no optimization, no randomness, no tolerance beyond $\varepsilon$.
\item[Warm start] The initial iterate $x^{(0)}$ of a step, taken to be the \emph{previous} 15-minute
step's certified solution. Not an approximation of the answer --- only a starting point whose
correctness is tested (see \emph{self-consistency}).
\item[Iterate] The value $x^{(r)}$ carried at round $r$ of the fixed-point loop. Each round produces
the next iterate by an exact solve; iterates are never averaged, relaxed, or damped.
\item[Self-consistency] The acceptance test: recomputing every agent's parameter point from the
\emph{solved} decision returns the same CRs that were assumed in order to compute it. Passing this
test certifies a GNE (\S\ref{sec:walkthrough}.6); nothing is accepted without it.
\item[Facet neighbour] A CR sharing a facet (a full-dimensional face) with a given CR.
\item[Iteration-free step] A step cleared with \textbf{exactly one decision broadcast per agent}. All
location, solving and refinement is local computation on the onboard maps and costs no communication,
however many rounds it takes.
\item[Fallback step] A step where the loop fails to certify and a warm-started distributed ADMM is run
instead --- the only part of the method in which agents exchange messages iteratively.
\end{description}
\end{keybox}

\begin{table}[H]\centering\small
\begin{tabular}{lll}
\toprule
Symbol & Meaning & Size / value\\
\midrule
\multicolumn{3}{l}{\emph{Indices and dimensions}}\\
$i,j$        & agent indices & $0,\dots,5$\\
$N$          & number of agents & $6$\\
$k$          & step index within the horizon & $0,\dots,H{-}1$\\
$H$          & receding horizon length & $4$ steps ($1$ hour)\\
$\Delta t$   & step duration & $0.25$ h\\
$t$          & real-time (RTM) step index within a day & $0,\dots,95$\\
$r$          & round index of the fixed-point loop & $1,\dots,4$\\
\midrule
\multicolumn{3}{l}{\emph{Decisions}}\\
$p_i$        & agent $i$'s grid import over the horizon & $\mathbb R^{4}$ (kW)\\
$cv_i$       & agent $i$'s renewable curtailment (renewable agents only) & $\mathbb R^{4}$ (kW)\\
$x_i$        & agent $i$'s decision block: $p_i$ (grid) or $[p_i;cv_i]$ (renewable) & $\mathbb R^{4}$ or $\mathbb R^{8}$\\
$x$          & stacked decision vector of all agents & $\mathbb R^{40}$\\
$x^\star$    & the certified equilibrium of a step & $\mathbb R^{40}$\\
\midrule
\multicolumn{3}{l}{\emph{Parameters}}\\
$\theta$     & \textbf{public} market data of the current step & $\mathbb R^{24}$\\
$D_i$        & H$_2$ still owed by agent $i$ over the horizon & kg\\
$\lambda_{\mathrm{RT}}$ & real-time price over the horizon & \$/MWh\\
$p_{\mathrm{DA},i}$ & agent $i$'s day-ahead award for the current hour & kW\\
$g_s,g_w$    & solar / wind availability at the shared farms & kW\\
$s_i$        & coalition aggregate excluding $i$: $\sum_{j\ne i}p_j$ & $\mathbb R^{4}$ (kW)\\
$t_i$        & agent $i$'s \textbf{private} parameter point $[\,s_i\,;\,\theta\,]$ & $\mathbb R^{28}$\\
$S$          & selector: extracts the first $H$ entries (the $p$-part) of $x_j$ & $\mathbb R^{4\times n_{x_j}}$\\
\midrule
\multicolumn{3}{l}{\emph{Physical and economic constants}}\\
$L_{\min},L_{\max}$ & PCC import band, binding on $\sum_ip_i$ at every step & $100,\ 900$ kW\\
$P_{\max,i}$ & agent $i$'s electrolyzer (converter) rating & kW\\
$\eta_i$     & agent $i$'s electrolyzer efficiency & kg/kWh\\
\midrule
\multicolumn{3}{l}{\emph{Maps and equilibrium system}}\\
$E_i,f_i$    & the polytope of one CR of agent $i$ & $E_it_i\le f_i$\\
$G_i,g_i$    & the affine law carried by that CR & $x_i=G_it_i+g_i$\\
$G_i^{s},G_i^{\theta}$ & $G_i$ split at column $H$, matching $t_i=[s_i;\theta]$ & --- \\
$r_i$        & index of the CR located for agent $i$ & --- \\
$\mathcal C$ & combination: one CR per agent, $(r_0,\dots,r_5)$ & --- \\
$M_x,M_p,M_1$& the equilibrium system $M_xx=M_p\theta+M_1$ & $40{\times}40$, $40{\times}24$, $\mathbb R^{40}$\\
$V_2$        & basis of $\ker M_x$ when $M_x$ is rank-deficient & --- \\
$y_2$        & coordinate on the GNE manifold, $x=x_p+V_2y_2$ & $\mathbb R^{40-\mathrm{rank}M_x}$\\
\midrule
\multicolumn{3}{l}{\emph{Cost}}\\
$Q,c,F$      & stacked cost: $\Phi(x)=\tfrac12x^\top Qx+(c+F\theta)^\top x$ & $Q\succ0$, $\lambda_{\min}=10^{-3}$\\
$\Phi$       & coalition potential; its minimizer over the feasible set is the v-GNE & --- \\
$G,w_0,W$    & joint feasible set of the coalition, $Gx\le w_0+W\theta$ & $142\times40$\\
\bottomrule
\end{tabular}
\caption{Nomenclature. Sizes are those of the case study of this report ($N{=}6$, $H{=}4$).}
\label{tab:nomen}
\end{table}

\section{Method overview and contributions}\label{sec:overview}
% ─────────────────────────────────────────────────────────────────────────────
This section states the problem, the method, and what is new, in full but without detail. The complete
derivation, line-by-line algorithm and two traced numerical examples are in
\S\ref{sec:walkthrough}; readers who want the mechanism should go there.

\subsection{The problem}
Six independently-owned green-hydrogen electrolyzers sit behind one point of common coupling (PCC) and
bid as a coalition into ERCOT. Every $15$ minutes the market publishes $\theta$ --- prices, renewable
availability, each agent's remaining H$_2$ obligation, the day-ahead awards --- and the coalition must
produce a dispatch $x^\star$ that is a \textbf{generalized Nash equilibrium}: no agent can lower its
own cost by deviating, and the coalition's total import obeys the physical band
$L_{\min}\le\sum_ip_i\le L_{\max}$ at every step.

The binding difficulty is not the equilibrium itself but the ownership model. \textbf{No agent
discloses its cost model}, so no party --- not even a hypothetical coordinator --- can assemble the
joint problem at runtime. Any admissible method must be built from quantities each agent can compute
alone, plus data that is already public. The only such coupling quantity is the coalition's total
import, which is already measured at the PCC meter.

\subsection{What FACET-RHG does}
Offline and once, each agent solves \emph{its own} mpQP and stores the result as a lookup table of
critical regions. Online, at each step, the coalition runs a \textbf{fixed-point iteration on the
combination}, summarized in Figure~\ref{fig:loop}.

\begin{figure}[H]\centering
\begin{tcolorbox}[colback=blue!2!white,colframe=blue!50!black,width=0.94\textwidth,
                  title={\bfseries One 15-minute step: warm start $\to$ locate $\to$ solve $\to$ certify}]
\small
\begin{enumerate}[leftmargin=1.9em,nosep,itemsep=3pt]
\item[\textbf{0.}] \textbf{Warm start.} Take $x^{(0)}$ to be the previous step's \emph{certified}
solution. Steps are $15$ minutes apart, so $\theta$ has moved little.
\item[\textbf{1.}] \textbf{Locate} (each agent, locally, in parallel). Agent $i$ forms its private
parameter point $t_i=[\,s_i\,;\,\theta\,]$, where $s_i=\sum_{j\ne i}p_j$ is read off the current
iterate, and finds the CR containing it by the containment test $E_it_i\le f_i$. Deterministic
arithmetic --- no optimization, no randomness. Six lookups give a combination $\mathcal C$.
\item[\textbf{2.}] \textbf{Solve exactly.} \emph{Given} $\mathcal C$, every agent's law is affine, so
the six coupled best responses collapse into one linear system $M_xx=M_p\theta+M_1$ whose exact
solution returns \emph{all six} agents' decisions at once. Where $M_x$ is rank-deficient --- the
coupling ceiling --- the GNE form a manifold and we select the \emph{variational} one by minimizing
the potential $\Phi$ over it.
\item[\textbf{3.}] \textbf{Certify.} Recompute each $t_i$ from the \emph{solved} $x$ and locate again.
If every agent returns the CR it was assumed to be in, the combination has reproduced itself: each
agent is provably best-responding to the aggregate that actually materializes, which \emph{is} the
definition of a GNE. Accept and stop. Otherwise adopt $x$ as the next iterate and return to 1.
\end{enumerate}
\tcblower
\small Steps 1--3 are \textbf{local computation on the onboard maps}. The inter-agent exchange is
\textbf{exactly one decision broadcast per agent per step}, regardless of how many rounds run.
\end{tcolorbox}
\caption{The online method. Nothing is accepted that has not passed step 3.}
\label{fig:loop}
\end{figure}

Two properties are worth stating plainly, because both are easy to mis-read:
\begin{itemize}[leftmargin=1.4em,nosep]
\item \textbf{The warm start is not an approximation of the answer.} It is only a starting point, and
its correctness is \emph{tested} in step 3. A poor warm start costs extra rounds of local arithmetic,
never a wrong dispatch. On 6~April, $82$ of $96$ steps certified on the first round.
\item \textbf{The method is sound, not complete.} Step 3 never accepts a wrong answer, but the loop is
not guaranteed to reach a fixed point. When it does not, the step falls back (\S\ref{sec:facet}); the
cost is a slower step, never a silently incorrect one.
\end{itemize}

\subsection{What is new here}
The architecture --- per-agent offline maps, a block linear system of affine best responses, membership
validation, one exchange per step --- is inherited from our IF-mpDiMPC work (Saini \emph{et al.},
\emph{Comput.\ Chem.\ Eng.}\ 200, 2025) and its ACC-2026 successor. \textbf{Point location is not new:
IF-mpDiMPC-V2 already performs it.} What is new is what the \emph{game} forces, none of which arises in
cooperative distributed MPC, where all controllers share one objective and the solution is unique:

\begin{enumerate}[leftmargin=1.7em,nosep,itemsep=3pt]
\item \textbf{Rank-deficient systems carry the answer instead of signalling an error.} IF-mpDiMPC
discards singular systems (\texttt{if abs(det(L)) < 1e-10: return None}) --- correct there, because a
unique solution means singularity implies a wrong combination. Here $M_x$ is genuinely rank-deficient
whenever the PCC band binds for several agents at once (rank $36$ of $40$ in the traced example of
\S\ref{sec:walkthrough}.8), the equilibria form a manifold, and the answer lies on it. We select the
variational GNE from that manifold, and re-impose feasibility on the selection when required
(Tier~1c).
\item \textbf{The iterate is refined rather than enumerated around.} IF-mpDiMPC-V2 and FACET-DiMPC
point-locate once and compensate by enumerating the product $\prod_i(1+|\mathcal N_i|)$ over each
agent's facet neighbours --- $985{,}608$ combinations at $t{=}0$ for this fleet. We instead re-locate
and iterate, solving a mean of $\mathbf{1.2}$ combinations per step. This matters because the locality
assumption underpinning the neighbour product fails in this game: when the combination changes,
$93.1\%$ (April) / $92.1\%$ (July) of the changes are more than three facet-hops away
(\S\ref{sec:walkthrough}.10).
\item \textbf{The coupling is sum-compressed.} IF-mpDiMPC's parameter vector carries every other
controller's full input trajectory, growing as $M\!\cdot\!N_p\!\cdot\!n_u$; ours carries four numbers,
because the PCC band is the only coupling. This is simultaneously the scalability argument and the
privacy argument --- the coalition total is already on the meter.
\end{enumerate}

\subsection{Headline results}
Over two independent ERCOT weeks (1--7 April and 7--13 July 2025; prices $-\$1$ to $\$3553$/MWh;
$1344$ real-time steps):
\begin{center}\small
\begin{tabular}{ll}
\toprule
Iteration-free steps & $\mathbf{99.3\%}$ \ ($9$ fallbacks of $1344$)\\
Communication on a map step & exactly $1$ decision broadcast per agent\\
Median real-time step & $\mathbf{125}$~ms against a $900$~s dispatch interval\\
Agreement with the centralized equilibrium & $\sim\!10^{-4}$~kW in the interior; $5.9$~kW worst case\\
Combinations ever visited & $700$ of $K^N\approx1.1\times10^{21}$ \ ($6.5\times10^{-19}$)\\
Weekly H$_2$ target & met at $140$--$141\%$\\
Renewable curtailment & $5.8$--$6.8\%$ of available generation\\
\bottomrule
\end{tabular}
\end{center}
\S\ref{sec:limit} reports honestly that the residual degeneracy at the coupling ceiling --- not the
fleet composition --- is what bounds the iteration-free rate, and \S\ref{sec:related} states which
comparisons in this report are not yet on an equal footing.

\section{Physical system and market}
A cluster of small electrolyzers sits behind one shared PCC. Each owner runs its own asset for its
own revenue and keeps its cost model private; only power decisions are shared. Two forces make this
a game rather than $N$ independent problems:
\begin{itemize}[nosep]
\item \textbf{Physical coupling.} The PCC caps the collective import and a market floor sets a
minimum offer, \emph{at every 15-min step $k$}:
\[
   L_{\min}=100\ \text{kW}\ \le\ \sum_{i=1}^{N} p_{i,k}\ \le\ L_{\max}=900\ \text{kW}.
\]
$L_{\min}$ is the ISO minimum offer size; $L_{\max}$ is the physical interconnection limit
($75\%$ of the $1200$~kW fleet nameplate). \textbf{This same band governs both the day-ahead and the
real-time stage} --- there is one physical wire.
\item \textbf{Private hydrogen contracts.} Each agent must produce a daily H$_2$ quantity; it may
time production freely but the daily total is a hard floor. That freedom, under a shared per-step
power limit, is what couples the agents and the horizon.
\end{itemize}
The market is ERCOT two-settlement: a day-ahead (DAM, hourly) stage that sets a financial anchor
$p^{\mathrm{DA}}$, and a real-time (RTM, 15-min SCED) stage that re-optimizes against live prices.
Electricity $\to$ H$_2$ conversion: $\text{kg}=\eta_i\,p_{\mathrm{elec},i}\,\Delta t$, $\Delta t=0.25$~h.

\section{The fleet (medium-scale)}
\begin{table}[H]\centering
\begin{tabular}{clcccccc}
\toprule
$i$ & Agent & type & $P_{\max,i}$ & farm cap & $\eta_i$ & $\gamma_i$ & break-even $a_i$\\
 & & & [kW] & [kW] & [kg/kWh] & & [\$/MWh]\\
\midrule
0 & PEM\_Elec  & grid & 250 & --  & 0.0200 & $5\!\times\!10^{-3}$ & 60.0\\
1 & ALK        & grid & 200 & --  & 0.0180 & $5\!\times\!10^{-3}$ & 54.0\\
2 & PEM\_PV    & PV   & 150 & 125 & 0.0210 & $5\!\times\!10^{-3}$ & 63.0\\
3 & PEM\_PV\_2 & PV   & 100 & 125 & 0.0190 & $5\!\times\!10^{-3}$ & 57.0\\
4 & PEM\_Wind  & wind & 275 & 250 & 0.0195 & $4\!\times\!10^{-3}$ & 58.5\\
5 & PEM\_Wind\_2 & wind & 225 & 250 & 0.0185 & $4\!\times\!10^{-3}$ & 55.5\\
\bottomrule
\end{tabular}
\caption{$N=6$, total electrolyzer nameplate $\sum_i P_{\max,i}=1200$~kW against
$L_{\max}=900$~kW, so the coupling ceiling binds. $r_{\mathrm{H2}}=\$3$/kg, $H=4$,
$\Delta t=0.25$~h. $P_{\max,i}$ is the \emph{electrolyzer} rating; ``farm cap'' is the
co-located renewable array, which is shared within a technology (both PV agents see the same
irradiance $g_h=\mathrm{CF}_h\cdot125$~kW, both wind agents the same $\mathrm{CF}_h\cdot250$~kW).
Break-even $a_i=r_{\mathrm{H2}}\eta_i\!\cdot\!10^3$: below $a_i$ buying grid power to make H$_2$
is profitable, above it the agent wants to stop. \textbf{All six agents are distinct in
$(P_{\max,i},\eta_i)$}, so all six $a_i$ are distinct and each agent requires its own offline
solve. Agents 3 and 5 are inverter-limited ($\text{farm}>P_{\max}$) and therefore curtail
$25$~kW at peak resource regardless of price.}
\end{table}

\begin{warnbox}{Distinct efficiencies remove one degeneracy but do not improve the iteration-free rate}
An earlier version of this study used two \emph{identical} PV agents and two identical wind agents.
Identical agents are ill-posed at the coupling ceiling in an obvious way: if agents $j$ and $k$ have
the same cost and constraints, any split of their combined power is equally optimal, so a continuum
of GNE exists. Making every $\eta_i$ distinct strictly orders the marginal H$_2$ values
$a_i=r_{\mathrm{H2}}\eta_i$, so ``who backs off first'' has a unique answer and \emph{that}
tie disappears by construction.

\textbf{Measured, this alone did not help.} Holding the clearing algorithm fixed, the
identical-pair fleet was $99.2\%$ iteration-free ($11/1344$ ADMM fallbacks) while the distinct
fleet was $97.2\%$ ($37/1344$). Permutation symmetry was never the \emph{binding} source of
degeneracy: per Hall \& Bemporad (Def.~2), $M_x$ loses rank whenever the \emph{shared coupling
constraint} is active for two or more agents at once --- which happens regardless of whether those
agents are identical. Removing the symmetry eliminated a secondary tie while the dominant one
remained, and the distinct fleet additionally brought richer maps and a new degenerate surface
(forced curtailment).

\textbf{Fixing the real cause recovers it, in two steps.} The dominant degeneracy is handled by the
feasibility-constrained variational selection of \S\ref{sec:facet} (Tier~1c), which addresses the
shared constraint rather than the fleet: with it alone the distinct fleet reached $98.8\%$
($16/1344$). A second, independent fix then closed the rest of the gap: the offline
facet-neighbour graph feeding Tiers 1b/1c was itself badly incomplete (\S\ref{sec:facet},
``\texttt{facet\_crossing}''); rebuilding it correctly brought the distinct fleet to
\textbf{$99.3\%$ ($9/1344$)} --- matching the identical-pair fleet's $99.2\%$. \emph{Fleet
composition was never the lever; equilibrium selection and a complete neighbour graph are.}
\end{warnbox}

% ─────────────────────────────────────────────────────────────────────────────
\section{Day-Ahead stage: distributed ADMM (model-private)}\label{sec:dam}
% ─────────────────────────────────────────────────────────────────────────────
The day-ahead schedule is solved once per day on the DAM \emph{forecast}
($\hat\lambda_h=\lambda^{\mathrm{DA}}_h+\mathcal N(0,16)$, renewable $\hat g_{i,h}$). It is
\textbf{distributed}: each agent $i$ solves only its own 24-hour problem, and the \emph{only}
exchanged quantity is the per-hour coalition import $\sum_i p_{i,h}$ (already public at the metered
point). No agent reveals $P_{\max,i}$, $\eta_i$, its renewable cap, or its daily target.

\paragraph{Per-agent day-ahead problem.} With $\Delta t_{\mathrm{DA}}=1$~h and daily target
$D^{\mathrm{day}}_i=0.55\,P_{\max,i}\eta_i\!\cdot\!24$, agent $i$ solves
\[
\min_{p_{i,h}\ge0,\;cv_{i,h}\ge0}\ \sum_{h=0}^{23}\Big[(\hat\lambda_h/1000-r_{\mathrm{H2}}\eta_i)\,
   \Delta t_{\mathrm{DA}}\,p_{i,h}+r_{\mathrm{H2}}\eta_i\,\Delta t_{\mathrm{DA}}\,cv_{i,h}
   +\underbrace{\tfrac12\gamma_{\mathrm{DA}}\,p_{i,h}^2}_{\text{regularizer}}\Big]
\]
\[
\begin{aligned}
&0\le p_{i,h}\le P_{\max,i},\quad 0\le cv_{i,h}\le \mathrm{cap}_i \qquad\text{(local box; $cv$ only renewables)}\\
&0\ \le\ p_{i,h}+\hat g_{i,h}-cv_{i,h}\ \le\ P_{\max,i} \qquad\text{(local renewable balance)}\\
&\textstyle\sum_{h=0}^{23}\eta_i\,(p_{i,h}+\hat g_{i,h}-cv_{i,h})\,\Delta t_{\mathrm{DA}}\ \ge\ D^{\mathrm{day}}_i
   \qquad\text{(hard daily H$_2$ floor)}
\end{aligned}
\]
Every term uses only agent $i$'s own data; the daily H$_2$ floor couples that agent's own 24 hours.
The \textbf{coupling band} $L_{\min}\le\sum_i p_{i,h}\le L_{\max}$ ($\forall h$) is the \emph{only}
cross-agent constraint, enforced by ADMM consensus/dual updates that see nothing but the per-hour
aggregate:
\[
\text{(x)}\ p_i^{\ell+1}=\arg\min\big[J_i(p_i)+\tfrac{\rho}{2}\|p_i-z_i^\ell\|^2\big],\quad
\text{(z)}\ z^{\ell+1}=\Pi_{[L_{\min},L_{\max}]}\!\big(\textstyle\sum_i p_i^{\ell+1}+\mu^\ell/\rho\big),\quad
\text{($\mu$)}\ \mu_i^{\ell+1}=\mu_i^{\ell}+\rho(p_i^{\ell+1}-z_i^{\ell+1}),
\]
with $\rho=0.1$; $\Pi$ is a closed-form clip of the per-hour aggregate onto the band. The award
$p^{\mathrm{DA}}_{i,h}$ is the RTM anchor; its cumulative-H$_2$ trajectory sets the RTM pacing target
$D_i(t)$. (\texttt{dam.solve\_dam\_admm}; consensus in \texttt{admm\_solver.py}.)

\begin{keybox}{The small $\gamma_{\mathrm{DA}}$ regularizer (why, and its trade-off)}
The day-ahead economics are \emph{linear} in each hourly buy: ``buy enough cheap power to make the
day's hydrogen.'' With no $\gamma_{\mathrm{DA}}$ the problem is a degenerate LP --- on a cheap day it
says \emph{buy the maximum every hour}, a knife-edge answer that ADMM oscillates on and cannot
settle. A \emph{tiny} quadratic penalty $\tfrac12\gamma_{\mathrm{DA}}p_{i,h}^2$
($\gamma_{\mathrm{DA}}=10^{-4}$) rounds the knife-edge: it spreads purchases smoothly, gives a
\emph{unique} schedule so ADMM converges, and changes total cost by $\le\!1.4\%$. A matching
$\varepsilon=10^{-3}$ regularizes renewable curtailment $cv$. Validation: a centralized solve of the
\emph{same} regularized problem (offline oracle, \texttt{solve\_dam\_centralized}) is matched to
$\le\!10$~kW on priced days; on ultra-flat days the day-ahead problem is genuinely degenerate so the
two solvers land on different equally-optimal schedules ($\le\!23$~kW apart, cost within $1\%$).
\end{keybox}

% ─────────────────────────────────────────────────────────────────────────────
\section{Real-Time stage: each agent's mpQP}
% ─────────────────────────────────────────────────────────────────────────────
Every RTM step, agent $i$ solves a strictly-convex mpQP over the horizon $k=0,\dots,3$. The decision
is grid import $p_{i,k}\ge0$ (renewable agents additionally choose curtailment $cv_{i,k}$). The cost
is a day-ahead deviation anchor $\tfrac12\gamma_i(p-p^{\mathrm{DA}}_i)^2$ (strict convexity,
$Q_i\succ0$) plus real-time energy cost net of H$_2$ revenue. \textbf{Each agent's cost and every
local constraint use only that agent's own parameters; the sole shared quantity is the aggregate
$\mathrm{s}_k=\sum_{j\ne i}p_{j,k}$.}

\subsection{Grid agents (PEM\_Elec $i{=}0$, ALK $i{=}1$)}
\begin{agentbox}{Grid agent $i$ --- decision $x_i=[p_{i,0..3}]\in\mathbb R^4$}
\textbf{Private parameters} $\theta_i=[\,\mathrm{s}_0..\mathrm{s}_3\mid D_i,\ \lambda_0..\lambda_3,\ p^{\mathrm{DA}}_i\,]\in\mathbb R^{10}$
\[
\min_{p_i\in\mathbb R^4}\ \sum_{k=0}^{3}\Big[\tfrac12\gamma_i\,(p_{i,k}-p^{\mathrm{DA}}_i)^2
   +\big(\lambda_k/1000-r_{\mathrm{H2}}\eta_i\big)\,\Delta t\;p_{i,k}\Big]
\]
\vspace{-8pt}
\[
\begin{aligned}
&0\le p_{i,k}\le P_{\max,i}, &&\text{(capacity box)}\\
&\textstyle\sum_{k}\eta_i\,p_{i,k}\,\Delta t\ \ge\ D_i, &&\text{(cumulative H$_2$ --- couples the 4 steps)}\\
&L_{\min}\ \le\ p_{i,k}+\mathrm{s}_k\ \le\ L_{\max}. &&\text{(shared per-step coupling)}
\end{aligned}
\]
\end{agentbox}

\subsection{Renewable agents (PEM\_PV $i{=}2,3$; PEM\_Wind $i{=}4,5$)}
Renewable output $g_{i,k}$ feeds the electrolyzer load $p_{\mathrm{elec},i,k}=p_{i,k}+g_{i,k}-cv_{i,k}$
(free renewable displaces grid import; curtailment $cv$ throws away H$_2$ at break-even).
\begin{agentbox}{Renewable agent $i$ --- decision $x_i=[p_{i,0..3},\,cv_{i,0..3}]\in\mathbb R^8$}
\textbf{Private parameters} $\theta_i=[\,\mathrm{s}_0..\mathrm{s}_3\mid D_i,\ \lambda_0..\lambda_3,\ p^{\mathrm{DA}}_i,\ g_0..g_3\,]\in\mathbb R^{14}$
\[
\min_{p_i,cv_i}\ \sum_{k=0}^{3}\Big[\tfrac12\gamma_i(p_{i,k}-p^{\mathrm{DA}}_i)^2+\tfrac12\varepsilon\,cv_{i,k}^2
  +(\lambda_k/1000-r_{\mathrm{H2}}\eta_i)\Delta t\,p_{i,k}+r_{\mathrm{H2}}\eta_i\,\Delta t\,cv_{i,k}\Big]
\]
\vspace{-8pt}
\[
\begin{aligned}
&0\le p_{i,k}\le P_{\max,i},\quad 0\le cv_{i,k}\le g_{i,k}\le\mathrm{cap}_i,\\
&0\ \le\ p_{\mathrm{elec},i,k}=p_{i,k}+g_{i,k}-cv_{i,k}\ \le\ P_{\max,i},\\
&\textstyle\sum_{k}\eta_i\,p_{\mathrm{elec},i,k}\,\Delta t\ \ge\ D_i,\qquad
L_{\min}\ \le\ p_{i,k}+\mathrm{s}_k\ \le\ L_{\max}\ \text{(grid import only)}.
\end{aligned}
\]
\end{agentbox}

\paragraph{Receding-horizon coupling (not a repeated static game).} The four steps couple only
through the cumulative-H$_2$ row. At an H$_2$-binding point (PEM, raise $\lambda_1$ by \$40/MWh):
$\Delta p=[+0.5,-1.5,+0.5,+0.5]$~kW with $\sum_k p_k:600\to600$ (H$_2$-conserved). Step~1 becomes
expensive so it drops; the others rise to keep the H$_2$ total --- a single-step map would give
$\Delta p=\mathbf0$.

\paragraph{Sum-compression and elimination.} Agent $i$ depends on the others only through the
per-step aggregate $\mathrm{s}_k=\sum_{j\ne i}p_{j,k}$ (4 numbers), not their private
$D_j,p^{\mathrm{DA}}_j,g_j$. Each agent's mpQP is solved offline (PPOPT) in its private $10$- or
$14$-dimensional $\theta_i$; the coupling variable is eliminated when the coalition combination is
assembled. Public parameter count is $4(\lambda)+6(p^{\mathrm{DA}})+6(D)+8(g)=24$, but no
24-dimensional problem is ever solved.

The daily pacing target is receding: $D_i(t)=\text{clip}\big((D^{\mathrm{day}}_i-h2^{\mathrm{inv}}_i)\,
H/\text{steps left in day},\,0,\,D^{\max}_i\big)$; only step $k=0$ is executed against realized price
and renewable, then the window advances.

% ─────────────────────────────────────────────────────────────────────────────
\section{Online FACET clearing of the RTM GNE}\label{sec:facet}
% ─────────────────────────────────────────────────────────────────────────────
Each agent's offline solution is an explicit piecewise-affine map: a set of \emph{critical regions}
(CRs), each a polytope $E_i\theta_i\le f_i$ carrying an affine law $x_i=G_i\theta_i+g_i$. A
\emph{combination} $\mathcal C=(r_0,\dots,r_5)$ picks one CR per agent; the equilibrium at combination
$\mathcal C$ solves a block-linear system $M_x x = M_p\theta+M_1$. FACET finds and solves the correct
combination online \textbf{without enumerating the $K^N\!\approx\!1.1\times10^{21}$ combinations}.

\subsection{How FACET works, step by step}\label{sec:walkthrough}
% ─────────────────────────────────────────────────────────────────────────────
This subsection builds the online method from nothing, in the order the code executes it, and then
runs two \emph{real} 15-minute steps end to end with the actual numbers. Every figure quoted is a
live trace of the shipped code on 6~April 2025 --- none of it is illustrative.

% ── 1 ────────────────────────────────────────────────────────────────────────
\subsubsection*{1. The question we must answer every 15 minutes}
At each real-time step the market hands us $\theta$ --- $24$ public numbers: the four upcoming RT
prices, the six agents' H$_2$ still owed, the six DAM awards, and the solar/wind available to each
farm. From $\theta$ alone we must produce
\[
  x^\star=\bigl(x_0,\dots,x_5\bigr),\qquad x_i=\text{agent }i\text{'s power over the next 4 steps},
\]
such that \textbf{(a)} no agent can reduce its own cost by unilaterally changing its own $x_i$, and
\textbf{(b)} the coalition respects the physical PCC band $L_{\min}\le\sum_i p_i\le L_{\max}$ at
every step. That is a \emph{generalized Nash equilibrium}. The word ``generalized'' is doing real
work: condition (b) couples the agents' \emph{feasible sets}, not merely their costs.

The hard constraint on any method is \S\ref{sec:facet}'s privacy model: \textbf{no agent reveals its
cost model}, so no party --- not even a hypothetical coordinator --- is able to write down the joint
problem at runtime. Whatever we do must be assembled out of things each agent can compute alone.

% ── 2 ────────────────────────────────────────────────────────────────────────
\subsubsection*{2. What each agent stores offline: a lookup table}
Before the day begins, each agent $i$ solved \emph{its own} mpQP once, and stored the answer as a
lookup table. The table has one row per \textbf{critical region} (CR). Each row holds
\begin{center}\small
\begin{tabular}{ll}
\toprule
a polytope & $E_i\,t_i\le f_i$ \quad --- ``does this row apply?''\\
an affine law & $x_i=G_i\,t_i+g_i$ \quad --- ``if it applies, here is my exact optimum''\\
\bottomrule
\end{tabular}
\end{center}
\emph{Using} the table is not an optimization. It is a lookup with two operations: find the row whose
polytope contains the query point $t_i$ (a containment test $E_it_i\le f_i$ --- pure arithmetic, no
solver, no randomness), then evaluate one affine formula. The tables are big but they are only
tables: $507$ rows for PEM\_Elec up to $12{,}782$ for PEM\_PV\_2, $35{,}348$ in total
(Table~\ref{tab:cr}).

We call the choice of one row per agent a \textbf{combination} $\mathcal C=(r_0,\dots,r_5)$. If we
knew the right combination, we would be finished --- that is the entire point of having the tables.

% ── 3 ────────────────────────────────────────────────────────────────────────
\subsubsection*{3. The one subtle part: what the query point contains}
What is $t_i$ --- the key we look the table up with? It is \emph{not} just $\theta$:
\[
  t_i \;=\; \Bigl[\;\underbrace{\textstyle\sum_{j\ne i} p_j}_{\text{4 numbers}}
  \;\Big|\;\underbrace{\theta}_{\text{24 numbers}}\;\Bigr]\;\in\;\mathbb R^{28}.
\]
Why does agent $i$'s own optimum depend on the others at all? \textbf{Only through the PCC band.}
Agent $i$ may draw whatever it likes provided $\sum_j p_j$ stays inside $[L_{\min},L_{\max}]$ --- so
the others constrain $i$ \emph{exclusively via their total}, never individually. That is why four
numbers suffice (one per horizon step), and it is also the whole privacy argument: the coalition
total is already visible at the PCC meter, so sharing it reveals nothing. Agent $i$ never sees any
$P_{\max,j}$, $\eta_j$, or cost.

This is also precisely what makes the problem a \emph{game} rather than six independent problems.

% ── 4 ────────────────────────────────────────────────────────────────────────
\subsubsection*{4. The circularity --- the only genuinely hard part}
\begin{warnbox}{The difficulty in one sentence}
To look up my row I need $\sum_{j\ne i}p_j$. But $\sum_{j\ne i}p_j$ is part of the answer I am trying
to compute. \textbf{The key depends on the value.}
\end{warnbox}
Everything else in FACET is machinery for breaking this circle without a central solver. If this
subsection is confusing, it is almost certainly \emph{this} point that is the source --- the rest is
bookkeeping around it.

% ── 5 ────────────────────────────────────────────────────────────────────────
\subsubsection*{5. Breaking the circle: warm start $\to$ locate $\to$ solve $\to$ certify}
The circle is broken by a \textbf{fixed-point iteration that carries its own proof of correctness}.

\begin{algorithm}[H]
\caption{One FACET clearing step (Tier 1). In: $\theta$, previous solution $x^{\mathrm{prev}}$,
previous combination $\mathcal C^{\mathrm{prev}}$. Out: a certified GNE $x^\star$.}
\label{alg:facet}
\begin{algorithmic}[1]
\State $x\gets x^{\mathrm{prev}}$;\quad $\mathcal C\gets\mathcal C^{\mathrm{prev}}$
       \Comment{\textbf{warm start}: previous step's certified solution}
\For{round $=1,\dots,4$}
  \For{each agent $i=0,\dots,5$}\Comment{local \& parallel; no messages}
    \State $t_i\gets\bigl[\;\sum_{j\ne i}p_j(x)\;\big|\;\theta\;\bigr]$
       \Comment{form my key from the current iterate}
    \State $r_i\gets\textsc{Locate}\bigl(t_i,\ \mathrm{hint}=\mathcal C_i\bigr)$
       \Comment{\textbf{locate} --- which row holds $t_i$?}
  \EndFor
  \State $\mathcal C'\gets(r_0,\dots,r_5)$
  \State $x'\gets\textsc{SolveExact}(\mathcal C',\theta)$
       \Comment{\textbf{solve} --- $M_xx=M_p\theta+M_1$, all six at once}
  \If{$\textsc{Locate}\bigl([\,\textstyle\sum_{j\ne i}p_j(x')\,|\,\theta\,]\bigr)=r_i$ for every $i$}
    \State \Return $x'$
       \Comment{\textbf{re-check passed} $\Rightarrow$ certified GNE; stop}
  \EndIf
  \State $x\gets x'$;\quad $\mathcal C\gets\mathcal C'$
       \Comment{someone crossed a boundary --- iterate}
\EndFor
\end{algorithmic}
\end{algorithm}

\medskip
Before any example, here is \textbf{every line of Algorithm~\ref{alg:facet}, with the exact operation
it performs}. Throughout, $x\in\mathbb R^{40}$ is the stacked decision vector and $\theta\in\mathbb
R^{24}$ the public data:
\begin{center}\small
\begin{tabular}{clcl}
\toprule
$i$ & agent & $n_{x_i}$ & block $x_i$ and its slot in $x$\\
\midrule
$0$ & PEM\_Elec    & $4$ & $p_i\in\mathbb R^{4}$ \hfill $x[0{:}4]$\\
$1$ & ALK          & $4$ & $p_i\in\mathbb R^{4}$ \hfill $x[4{:}8]$\\
$2,3$ & PEM\_PV, PEM\_PV\_2     & $8$ & $[\,p_i\,;\,cv_i\,]\in\mathbb R^{8}$ \hfill $x[8{:}16]$, $x[16{:}24]$\\
$4,5$ & PEM\_Wind, PEM\_Wind\_2 & $8$ & $[\,p_i\,;\,cv_i\,]\in\mathbb R^{8}$ \hfill $x[24{:}32]$, $x[32{:}40]$\\
\midrule
\multicolumn{4}{l}{$n_x^{\mathrm{tot}}=2\cdot4+4\cdot8=\mathbf{40}$,\qquad $n_\theta=\mathbf{24}$,\qquad $H=4$}\\
\bottomrule
\end{tabular}
\end{center}

\paragraph{Line 1 --- the warm start ($x\gets x^{\mathrm{prev}}$, $\mathcal C\gets\mathcal C^{\mathrm{prev}}$).}
We begin from the previous 15-minute step's certified answer and the rows it used. \emph{At the very
first step of a day there is no previous step}: the code seeds $x$ with the DAM award,
$p_i(k)\gets p_{\mathrm{DA},i}$ for every $k$, and sets $\mathcal C\gets\varnothing$ --- so line 5 has
no hint and performs a full scan exactly once per day. Every later step inherits the previous one.

\paragraph{Line 2 --- the round loop.}
At most 4 rounds. One point is essential and easy to miss: \textbf{$\theta$ is fixed for the whole
loop.} It is the data of \emph{this} step and does not change between rounds. The \emph{only} thing
that changes from round to round is the iterate $x$, and therefore the aggregates $s_i$ built from it.

\paragraph{Line 3 --- the per-agent loop.}
Every agent executes lines 4--5 on its own hardware, in parallel, using only its own table. No
messages are exchanged inside the loop.

\paragraph{Line 4 --- form the lookup key, $t_i\gets[\,\sum_{j\ne i}p_j(x)\;|\;\theta\,]$.}
Two concatenated pieces. First the aggregate of everyone else, read off the \emph{current iterate} $x$:
\[
  s_i \;=\; \sum_{j\ne i}p_j\;=\;\sum_{j\ne i} S\,x_j\;\in\;\mathbb R^{4},
  \qquad S x_j := \text{first } H{=}4 \text{ entries of } x_j\ (\text{the grid-import part}),
\]
(for renewable agents $S$ discards the curtailment block $cv_j$ --- only imported grid power crosses
the PCC). Then $\theta$ is appended:
\[
  t_i=\bigl[\,s_i\;;\;\theta\,\bigr]\in\mathbb R^{4+24}=\mathbb R^{28}.
\]
\textbf{Which $\theta$ is this?} \emph{This step's} $\theta$ --- rebuilt from scratch at every RTM
step out of the current market data: the new RT prices $\lambda_{\mathrm{RT}}$, the freshly updated
H$_2$ still owed $D_i$, the current solar/wind availability $g$, and (on the hour) the new DAM award
$p_{\mathrm{DA}}$. It is \emph{public} and therefore identical for all six agents; only the first four
entries of $t_i$ differ between agents. And per line 2, it is constant across rounds.

\paragraph{Line 5 --- locate, $r_i\gets\mathrm{Locate}(t_i,\ \mathrm{hint}=\mathcal C_i)$.}
The question asked is only: \emph{which row $r$ of my table satisfies $E_r\,t_i\le f_r$?} The
procedure is
\begin{enumerate}[leftmargin=1.8em,nosep,label=(\alph*)]
\item test the hint row: if $\max(E_{\mathcal C_i}t_i-f_{\mathcal C_i})\le\varepsilon$, return $\mathcal C_i$;
\item else test the hint's facet neighbours, return the first that contains $t_i$;
\item else scan all rows with one batched matrix product,
      $\mathrm{viol}=\max_{\text{rows}}\bigl(\mathbf E\,t_i-\mathbf f\bigr)$, and return the first
      index with $\mathrm{viol}\le\varepsilon$.
\end{enumerate}

\begin{keybox}{What ``hint'' means on line 5 --- and why it is safe}
The \textbf{hint} is simply \emph{the row this agent used last time}: on round 1 it is the row from
the previous 15-minute step; on rounds $2,3,4$ it is the row found in the previous round. It is
carried in $\mathcal C$.

The hint is \textbf{purely a speed device --- it cannot change the answer, only the time taken to find
it.} Whenever $t_i$ lies strictly inside a region, that region is the unique one containing it, so
(a), (b) and (c) all return the same row. If $t_i$ lands exactly on a shared facet, two rows contain
it and (a)--(c) may return different indices --- but the piecewise-affine map is \emph{continuous}, so
both rows' affine laws produce the identical $x_i$. Either answer is correct.

This is why the warm start can never corrupt the result: a bad hint costs a few wasted containment
tests, never a wrong equilibrium. Measured: the hint alone answers $86.1\%$ (April) / $77.3\%$ (July)
of lookups --- one test against $19$--$49$ inequality rows (\S9).
\end{keybox}

\paragraph{Line 7 --- assemble the combination, $\mathcal C'\gets(r_0,\dots,r_5)$.}
Pure bookkeeping: collect the six located row indices into one tuple.

\paragraph{Line 8 --- \texttt{SolveExact}: the line where the coupling is resolved.}
This is the heart of the method, so it is worth deriving rather than asserting. Row $r_i$ of agent
$i$'s table states its \emph{exact} optimum as an affine function of its key:
\[
  x_i \;=\; G_i\,t_i + g_i .
\]
Split $G_i$ at column $H{=}4$, matching the two pieces of $t_i$:
\[
  G_i=\bigl[\,\underbrace{G_i^{s}}_{n_{x_i}\times4}\;\big|\;\underbrace{G_i^{\theta}}_{n_{x_i}\times24}\,\bigr]
  \qquad\Longrightarrow\qquad
  x_i \;=\; G_i^{s}\,s_i \;+\; G_i^{\theta}\,\theta \;+\; g_i .
\]
Now substitute what $s_i$ actually is --- a linear function of the \emph{other agents' unknowns},
$s_i=\sum_{j\ne i}S\,x_j$ --- and move it to the left:
\begin{equation}
  \boxed{\;x_i \;-\; G_i^{s}\!\!\sum_{j\ne i}S\,x_j \;=\; G_i^{\theta}\,\theta + g_i\;}
  \qquad\text{one block row, for each } i=0,\dots,5 .
  \label{eq:blockrow}
\end{equation}
Equation~\eqref{eq:blockrow} is the essential step: the circularity of \S4 has become an ordinary
\emph{linear} coupling, because inside a known combination every agent's response is affine. Stacking
the six block rows gives
\[
  M_x\,x \;=\; M_p\,\theta + M_1,
\]
\[
  \underbrace{\bigl[M_x\bigr]_{ii}=I_{n_{x_i}}}_{\text{``what I do with my own }x_i\text{''}},\qquad
  \underbrace{\bigl[M_x\bigr]_{ij}=-\,G_i^{s}S\ \ (j\ne i)}_{\text{``how }j\text{'s import pulls me''}},\qquad
  \bigl[M_p\bigr]_{i}=G_i^{\theta},\qquad
  \bigl[M_1\bigr]_{i}=g_i,
\]
with $M_x\in\mathbb R^{40\times40}$, $M_p\in\mathbb R^{40\times24}$, $M_1\in\mathbb R^{40}$. The
$-G_i^sS$ block sits in the first four columns of agent $j$'s slot --- the only columns $i$ can
\emph{see}. Read structurally, $M_x=I-(\text{coupling})$. \textbf{Six mutually-dependent best
responses have collapsed into one linear solve that returns all six agents' decisions at once} --- no
iteration, no message, no tolerance.

\smallskip
\emph{Generic case ($M_x$ invertible).} $x'=M_x^{-1}(M_p\theta+M_1)$: the unique GNE of this
combination. Example~A is this case ($40\times40$, rank $40$).

\smallskip
\emph{Degenerate case ($M_x$ singular).} Example~B is this case (rank $36$). It happens when the PCC
band binds for several agents simultaneously: their laws all reduce to ``I take whatever the band
leaves me'', so the block rows become linearly dependent. The solutions then form a
$\bigl(40-\mathrm{rank}\,M_x\bigr)$-dimensional \emph{manifold}, obtained from the SVD of $M_x$:
\[
  x = x_p + V_2\,y_2,\qquad V_2=\text{null-space basis of }M_x,\qquad y_2\in\mathbb R^{40-\mathrm{rank}M_x}.
\]
Every point on it is a GNE. We select the \textbf{variational} one --- the social optimum --- by
minimizing the coalition potential $\Phi$ over the manifold:
\[
  y_2^\star=\arg\min_{y_2}\ \Phi\bigl(x_p+V_2y_2\bigr),\qquad
  \Phi(x)=\tfrac12x^\top Qx+q^\top x,\quad q=c+F\theta,
\]
which is a tiny unconstrained QP of dimension $1$--$4$, solved from its normal equations
\[
  \bigl(V_2^\top QV_2\bigr)\,y_2 \;=\; -\,V_2^\top\bigl(Qx_p+q\bigr).
\]
This is why a plain \texttt{numpy.linalg.solve} is not enough: on Example~B it would simply fail.
(When even this selection leaves the feasible set, Tier~1c re-imposes the joint feasibility rows on
the same minimization --- see the keybox in \S\ref{sec:facet}.)

\paragraph{Line 9 --- the re-check (the acceptance test).}
Take the freshly computed $x'$, recompute each agent's aggregate \emph{from it},
$s_i'=\sum_{j\ne i}S\,x_j'$, rebuild $t_i'=[\,s_i'\,;\,\theta\,]$ --- note $\theta$ is unchanged, only
the aggregate is new --- and locate once more. Accept \textbf{only} if all six agents return the same
rows $r_i$ they were assumed to be in. \S6 explains why passing this test is a proof.

\paragraph{Line 10 --- return.} The combination reproduced itself: $x'$ is a certified GNE. Stop.

\paragraph{Line 12 --- iterate.} At least one agent's key crossed a facet, so the affine laws used to
build $x'$ were the wrong ones. Adopt $x'$ as the next iterate and repeat. In practice this
converges immediately: on 6~April, $82$ of $96$ steps passed on round 1, $12$ needed 2, one needed 3,
one needed 4.

\begin{warnbox}{Honesty: the loop is not guaranteed to converge}
Nothing above proves the fixed point is reached, and it is not always. That is exactly why the round
count is capped at 4 and why Tiers~1b/1c/2 exist. What \emph{is} guaranteed is soundness, not
completeness: line 9 never accepts a wrong answer, so a failure to converge costs a fallback, never a
silently incorrect dispatch. Over both weeks the loop closes on $99.3\%$ of steps.
\end{warnbox}

% ── 6 ────────────────────────────────────────────────────────────────────────
\subsubsection*{6. Why self-consistency certifies a generalized Nash equilibrium}
Suppose the loop returns $x'$. The test that let it return says: \emph{when every agent recomputes
its key from the aggregate that $x'$ actually produces, it lands in the very row that was used to
compute $x'$.} Since that row's affine law \emph{is} agent $i$'s exact optimum for that key, each
$x_i'$ is agent $i$'s exact best response to the aggregate that actually materializes. That is the
definition of a generalized Nash equilibrium --- so a self-consistent combination is not
``close enough'', it \textbf{is} the equilibrium.

This also explains why nothing may be accepted without the test. A combination that has not
reproduced itself is a combination whose affine laws were evaluated at a key that never occurred;
its output is meaningless, however plausible the numbers look. Example B below is exactly this trap.

% ── 7 ────────────────────────────────────────────────────────────────────────
\subsubsection*{7. Example A --- an ordinary step (\textit{t}=1, 00:15): one round}
Market data: prices $\lambda_{\mathrm{RT}}=[11.84,\,11.28,\,13.63,\,13.58]$~\$/MWh (cheap night), no
sun ($g_s=0$), wind $g_w=168.7$~kW, DAM awards $p_{\mathrm{DA}}=[250,200,150,100,75,25]$~kW.
Warm start: last step's $\Sigma p=805.67$~kW at combination $(369,325,4001,9222,8,14)$.

\textbf{Locate.} Each agent forms its own key and tests rows:
\begin{center}\small
\begin{tabular}{clcccc}
\toprule
$i$ & agent & $\sum_{j\ne i}p_j$ (kW) & row found & \#\,ineq. & $\max(E_it_i-f_i)$\\
\midrule
0 & PEM\_Elec   & $555.67$ & $369$ / $507$    & $19$ & $-2.6\times10^{-8}$\\
1 & ALK         & $605.67$ & $325$ / $507$    & $19$ & $+8.5\times10^{-7}$\\
2 & PEM\_PV     & $655.67$ & $4001$ / $5613$  & $39$ & $+9.5\times10^{-8}$\\
3 & PEM\_PV\_2  & $705.67$ & $9222$ / $12782$ & $44$ & $-2.9\times10^{-7}$\\
4 & PEM\_Wind   & $727.74$ & $8$ / $5613$     & $49$ & $-3.7\times10^{-1}$\\
5 & PEM\_Wind\_2& $777.93$ & $14$ / $10326$   & $45$ & $-4.2\times10^{-1}$\\
\bottomrule
\end{tabular}
\end{center}
The last column is the containment test: $\le0$ means the key sits inside that polytope. All six are
inside, giving $\mathcal C'=(369,325,4001,9222,8,14)$.

\textbf{Solve.} $M_x$ is $40\times40$ of \textbf{rank 40 --- full rank, so a unique GNE}:
\[
  p_i(k{=}0)=[250.00,\;200.00,\;150.00,\;100.00,\;77.92,\;27.73]\ \text{kW},\qquad
  \Sigma p=805.64\ \text{kW}\in[100,900].
\]

\textbf{Re-check.} Re-locating at this $x'$ returns $(369,325,4001,9222,8,14)$ --- identical.
Certified in one round. Against the offline centralized oracle,
$\lVert x-x_{\mathrm{cent}}\rVert_\infty=9.0\times10^{-8}$~kW. (Sanity: at $\$12$/MWh power is cheap,
so the four grid/PV agents simply sit at their DAM awards.)

% ── 8 ────────────────────────────────────────────────────────────────────────
\subsubsection*{8. Example B --- a boundary crossing (\textit{t}=20, 05:00): two rounds}
This is the case worth studying, and it shows the re-check earning its place. Prices have risen to
$\lambda_{\mathrm{RT}}\approx26.9$~\$/MWh and wind has fallen to $g_w=110.4$~kW. Warm start:
$\Sigma p=892.42$~kW at combination $(506,506,9,12,10,1)$ --- already within $7.6$~kW of the
$L_{\max}=900$~kW ceiling.

\textbf{Round 1 --- locate.} From the warm start: $\mathcal C'=(366,314,9,12,10,1)$. Note agents 0
and 1 land in \emph{different} rows than the warm combination: $\theta$ moved since the last step, so
their keys moved too.

\textbf{Round 1 --- solve.} $M_x$ is $40\times40$ of \textbf{rank 36 --- rank-deficient, nullity 4}.
The coupling binds for several agents at once, so there is no longer a single equilibrium but a
\textbf{4-dimensional manifold} of them (this is the degeneracy of \S\ref{sec:limit}; a plain
\texttt{solve} would fail outright here). The potential-minimizing selection picks the social optimum
from that manifold:
\[
  p_i(k{=}0)=[212.71,\;153.38,\;141.03,\;90.73,\;176.07,\;126.07],\qquad
  \Sigma p=\mathbf{900.00}\ \text{kW}=L_{\max}\ \text{exactly}.
\]

\textbf{Round 1 --- re-check: FAILS.} Re-locating returns $(506,506,9,12,10,1)$: agents 0 and 1 are
\emph{not} in the rows we assumed. Rejected.

\textbf{Round 2.} Iterate with the new $x$. Locate $\to(506,506,9,12,10,1)$, whose $M_x$ has
\textbf{rank 40 --- full rank, unique GNE}:
\[
  p_i(k{=}0)=[211.33,\;152.00,\;141.03,\;90.73,\;176.07,\;126.07],\qquad
  \Sigma p=897.23\ \text{kW}.
\]
Re-check returns $(506,506,9,12,10,1)$ --- identical. Certified.
$\lVert x-x_{\mathrm{cent}}\rVert_\infty=4.4\times10^{-9}$~kW.

\begin{keybox}{What Example B teaches}
Round 1's answer \emph{looked} perfect: jointly feasible, ceiling respected, $\Sigma p$ sitting
exactly at $L_{\max}=900$~kW. It was nevertheless \textbf{wrong}, because pushing the aggregate to
$900$~kW moved agents 0 and 1 out of the very rows whose affine laws had been used to compute it ---
the laws were evaluated at a key that never occurred. Only the re-check exposes this. Cost of the
whole episode: two lookups and two $40\times40$ linear solves, all local, and still
\textbf{one broadcast for the step}.
\end{keybox}

% ── 9 ────────────────────────────────────────────────────────────────────────
\subsubsection*{9. ``If I already know last step's combination, why locate at all?''}
A fair objection, and the answer is measured rather than argued. We \emph{do} reuse it --- but not as
the answer, only as the \textbf{hint}: \textsc{Locate} tests the hinted row \emph{first} and returns
immediately if the key is still inside it. So when nothing has changed, ``locating again'' costs
\textbf{exactly one containment test} ($19$--$49$ inequality rows), not a scan of $12{,}782$ rows.

But the previous combination is the answer for the \emph{previous} $\theta$. Between steps the prices,
the wind, the H$_2$ still owed and (on the hour) the DAM award all move, so every key moves and may
cross a facet. Measured over two full days:

\begin{center}\small
\begin{tabular}{lcc}
\toprule
 & 6 April & 9 July\\
\midrule
Final combination \textbf{unchanged} from previous step & $66/95$ \ ($69.5\%$) & $57/95$ \ ($60.0\%$)\\
Final combination \textbf{changed}                      & $29/95$ \ ($30.5\%$) & $38/95$ \ ($40.0\%$)\\
\textbf{Unverified reuse} (skip locating) would \textbf{fail}& $\mathbf{29/95}$ \ $(\mathbf{30.5\%})$ & $\mathbf{28/95}$ \ $(\mathbf{29.5\%})$\\
\midrule
\multicolumn{3}{l}{\emph{Where the lookup found the row (round 1; $576$ agent-lookups/day):}}\\
\quad hint hit --- one test        & $496$ \ ($86.1\%$) & $445$ \ ($77.3\%$)\\
\quad facet-neighbour of the hint  & $6$ \ ($1.0\%$)    & $38$ \ ($6.6\%$)\\
\quad full vectorized scan         & $74$ \ ($12.8\%$)  & $93$ \ ($16.1\%$)\\
\bottomrule
\end{tabular}
\end{center}

So the previous row remains valid $\sim\!86\%$ of the time --- and the one cheap test is exactly what tells us
\emph{which} case we are in. Simply assuming last step's combination fails on $\sim\!30\%$ of steps.
Note also that ``locate again'' and ``verify the previous combination is still valid'' are the
\emph{same operation}: there is no cheaper check being skipped.

One further reading of that table matters. Of the $80$ April lookups that missed the hint, only $6$
landed on a facet \emph{neighbour} --- the other $74$ ($92.5\%$) needed a full scan. When an agent's
row changes, the new row is usually \textbf{not adjacent} to the old one.

% ── 10 ───────────────────────────────────────────────────────────────────────
\subsubsection*{10. ``Why not just walk the neighbour graph, as in FACET-DiMPC?''}
Because of what \S9's last paragraph and Example B both show. A neighbour walk assumes the
combination moves \emph{locally}: the parameter shifts slightly, you cross one facet, the new
combination is one hop away. That holds when the solution is unique and varies continuously --- the
cooperative DiMPC setting the walk was designed for. At a rank-deficient ceiling it does not: the
selection can jump across the manifold, and consecutive combinations end up graph-far apart. In
Example B agents 0 and 1 jumped from rows $(366,314)$ to $(506,506)$ in a single step. Point location
is indifferent to graph distance --- it goes straight to the right row in one containment test.

Head-to-head on 6~April, same seed, neither method allowed an ADMM fallback:
\begin{center}\small
\begin{tabular}{lcc}
\toprule
& facet-neighbour walk only & point location (Tier 1)\\
\midrule
Steps solved & $77/95$ \ ($81.1\%$) & $\mathbf{94/95}$ \ $(\mathbf{98.9\%})$\\
Combinations solved per step & mean $776$, max $4000$ & mean $\mathbf{1.2}$, max $4$\\
Worst $\lVert x-x_{\mathrm{cent}}\rVert_\infty$ on solved steps & $28.2$~kW & $\mathbf{9.2\times10^{-5}}$~kW\\
\bottomrule
\end{tabular}
\end{center}
The last row matters as much as the first: when the walk \emph{does} land on a combination, it
sometimes lands on a \emph{different equilibrium} $28$~kW from the social optimum. The walk is
retained as a backup (Tier~1b) because it does rescue occasional steps point location misses, and its
graph is load-bearing for the min-potential refinement that runs on \emph{every} step --- but it is
not the primitive that does the work in this game. (One day, $95$ steps, one seed; the walk arm is our
re-implementation of the ACC-2026 scheme, capped at 3 hops / $4000$ combinations.)

% ── 11 ───────────────────────────────────────────────────────────────────────
\subsubsection*{11. What happens when the loop does not close}
On 6~April the loop above certified \textbf{82 of 96 steps in a single round}; $12$ took two rounds,
one took three, one took four. Across both full weeks it clears $99.3\%$ of steps. The remaining
$0.7\%$ sit at genuinely degenerate ceiling vertices, and are handled by the backups summarized
next: the facet-neighbour walk (Tier~1b), the feasibility-constrained variational solve (Tier~1c),
and --- last and rarest --- a warm-started ADMM fallback (Tier~2), the only part of the method that
iterates between agents. The keybox below is the compressed statement of all of it.

\begin{keybox}{The FACET method (what runs every 15-minute step)}
\begin{enumerate}[leftmargin=1.6em,nosep]
\item \textbf{Per-agent point-location (warm-started).} Each agent locates its own CR at
$[\mathrm{s}_i;\theta]$ using the \emph{previous} step's CR as a hint: test that CR and its
facet-neighbours first (\,mean $5.9$--$8.9$ neighbours/CR, see below\,); only on a miss does a
vectorized scan over that agent's own CRs run. This is per-agent (thousands of CRs), never over
combinations.
\item \textbf{Self-consistent coalition combination.} Assemble the located CRs into a combination;
re-locate each agent given the others' aggregate and iterate to a fixed point --- a genuine GNE
(each agent best-responding to the shared aggregate).
\item \textbf{Variational-GNE solve.} Solve $M_x x=M_p\theta+M_1$ by the \emph{potential-minimizing}
selection: when the coupling binds for $\ge2$ agents $M_x$ is rank-deficient (a manifold of
equilibria); we pick the one minimizing the coalition potential $\sum_i J_i$ = the variational GNE =
social optimum (matches the centralized oracle). A plain solve would fail on that singular system.
\item \textbf{Facet-neighbour walk (backup).} If the located combination is not yet a strict
equilibrium, walk the combination's facet-neighbour graph (1st/2nd/3rd neighbour) for the closest
valid one; a 1-hop min-potential refinement locks the variational combination. The graph itself is
built offline by \emph{physically crossing every facet} (\texttt{facet\_crossing}, below) rather than
matching hyperplane coefficients --- see the keybox after Table~\ref{tab:cr}.
\item \textbf{Feasibility-constrained variational solve (Tier 1c).} If steps 1--4 accept nothing,
re-solve the best located combination with the joint feasibility rows \emph{imposed on the
selection}: pick $y_2$ by $\min_{y_2}\Phi(x_p+V_2y_2)$ s.t.\ $G(x_p+V_2y_2)\le w_0+W\theta$, the
role of the paper's Eq.~(9b) in this sum-compressed formulation. This is a tiny QP --- its dimension
is the nullity of $M_x$, typically $1$--$4$ --- and purely local, so the step still costs one
broadcast. Recovers the v-GNE to $\sim\!10^{-3}$~kW at a $\sim\!10^{-13}$ feasibility residual;
measured on the (then-incomplete) hash-based neighbour graph it rescued $21$ of the $37$ would-be
fallbacks ($97.2\%\!\to\!98.8\%$). Combined with the subsequent \texttt{facet\_crossing} graph fix
below, the two together reach $99.3\%$ ($9/1344$).
\item \textbf{ADMM fallback (rare; the only iteration).} If steps 1--5 fail --- only at genuinely
degenerate ceiling vertices --- fall back to a warm-started distributed ADMM solve. Its result is
used only if it converges; a non-converged result never becomes a warm start.
\end{enumerate}
All of steps 1--5 are \emph{local computation on the onboard maps}: the inter-agent exchange is
\textbf{exactly one decision broadcast per agent per step}, regardless of how many local checks run.
\end{keybox}

\begin{warnbox}{Why the selection is constrained only in Tier 1c, never in steps 3--4}
The constraint in Tier~1c must \textbf{not} be folded into the ordinary variational solve of step 3.
In steps 3--4 the \emph{unconstrained} selection's infeasibility is a load-bearing signal: it is how
the search learns that a candidate combination is wrong and keeps walking. Constraining the solve
everywhere manufactures a feasible-looking point at wrong combinations, so the walk stops early on
them --- measured, that raised fallbacks on the four worst days from $7/3/4/0$ to $15/35/29/17$ and
\texttt{map=cent} from $\sim\!3$~kW to $12$--$22$~kW. Constraining with the CR-membership rows
$E_i\theta_i\le f_i$ instead of the joint feasibility rows fails the same way, by forcing the point
onto the combination's CR corner. Tier~1c is safe precisely because it runs \emph{only} on steps that
would otherwise iterate, where the search has already exhausted its candidates.
\end{warnbox}

\begin{keybox}{Privacy model (no centralized optimizer at either stage)}
Agents share \textbf{only their bidding decisions, never their cost models} --- in both the DAM
(\S\ref{sec:dam}) and the RTM (\S\ref{sec:facet}). Consequently no party can solve the joint problem
at runtime. A centralized QP is used \textbf{only offline} to certify accuracy (map $=$ centralized)
and to bound an iterative solver's round count; it is never in the loop.
\end{keybox}

\subsection{The multiparametric maps (built once, offline)}
Each agent's RTM mpQP is solved \emph{once, offline} in its private parameter space $\theta_i$ by a
combinatorial multiparametric-QP solver (PPOPT). The result is the explicit piecewise-affine map:
a set of critical regions (CRs), each a polytope $E_i\theta_i\le f_i$ carrying an affine control law
$x_i=G_i\theta_i+g_i$. Because every agent differs in $(P_{\max,i},\eta_i)$, \textbf{six distinct
solves} are required, one per agent, each then expanded into its own slot of the public
parameter vector. Crucially, the private
$\theta_i$ box spans the \emph{full} operating range of every parameter --- price
$\lambda\in[-50,300]$, day-ahead anchor $p^{\mathrm{DA}}_i\in[0,P_{\max,i}]$, demand
$D_i\in[0,D^{\max}_i]$, renewable $g$ --- so the \emph{same} maps are valid for \emph{every} day and
\emph{every} DAM anchor. \textbf{The RTM maps are therefore built exactly once, never rebuilt}; only
the day-ahead ADMM runs daily.

\begin{table}[H]\centering
\begin{tabular}{llcclr}
\toprule
Agent & type & $\dim\theta_i$ & stack/farm [kW] & decision $x_i$ & CRs\\
\midrule
PEM\_Elec   & grid      & 10 & $250/$--    & $[p_{0..3}]\in\mathbb R^4$           & 507\\
ALK         & grid      & 10 & $200/$--    & $[p_{0..3}]\in\mathbb R^4$           & 507\\
PEM\_PV     & renewable & 14 & $150/125$   & $[p_{0..3},cv_{0..3}]\in\mathbb R^8$ & 5{,}613\\
PEM\_PV\_2  & renewable & 14 & $100/125$   & $[p_{0..3},cv_{0..3}]\in\mathbb R^8$ & 12{,}782\\
PEM\_Wind   & renewable & 14 & $275/250$   & $[p_{0..3},cv_{0..3}]\in\mathbb R^8$ & 5{,}613\\
PEM\_Wind\_2& renewable & 14 & $225/250$   & $[p_{0..3},cv_{0..3}]\in\mathbb R^8$ & 10{,}326\\
\midrule
\textbf{Total} & & & & & \textbf{35{,}348} (6 distinct solves)\\
\bottomrule
\end{tabular}
\caption{Offline multiparametric maps: per-agent critical-region counts. Each agent solves its own
$10$- or $14$-dimensional mpQP; $\theta_i$ is compressed by the coupling ``sum'' variable so no
$24$-dimensional problem is ever solved. One-time offline cost $\approx\!26$~min for all six solves
(12-core, exhaustive combinatorial mpQP); the facet-neighbour graph (\texttt{facet\_crossing},
\S\ref{sec:facet}) adds $\approx\!16$~min, giving $285{,}887$ neighbour pairs with \textbf{zero}
CRs left without a recorded neighbour.
\textbf{Generated once} and reused for both entire weeks.
\newline\textbf{Map size tracks the stack/farm ratio, not capacity.} The two agents whose stack
exceeds their farm ($150/125$, $275/250$) yield $5{,}613$ CRs each: once the array saturates, the
converter ceiling $p+g-cv\le P_{\max}$ can no longer bind, and a whole family of active sets
vanishes. The two inverter-limited agents ($100/125$, $225/250$) instead yield $12{,}782$ and
$10{,}326$: their stack ceiling and curtailment bounds bind over a much wider slice of parameter
space. The effects nearly cancel, so moving from $4$ shared solves to $6$ distinct ones grows the
total map by only $9\%$ ($32{,}402\to35{,}348$).}
\label{tab:cr}
\end{table}

\begin{keybox}{Finding facet neighbours: why hyperplane-matching fails at this CR count}
Two standard methods identify facet-adjacent CRs by matching normalized hyperplane coefficients
between region pairs: a hash bucket (fast, used in an earlier build of these maps) or, per the
FACET-DiMPC paper (ACC 2026, Eq.~14) this method builds on, an exhaustive pairwise LP scan. Both
require the shared boundary to be represented \emph{algebraically identically} by both regions'
constraint rows. Ground-truthed against physically crossing each facet (below): the hash method
recovered only $\sim\!12\%$ of the true neighbour count on this fleet's maps, leaving most CRs with
far too few recorded neighbours --- despite CRs geometrically tiling the parameter space (every
full-dimensional CR must have $\ge1$ neighbour, except at the $\theta$-box boundary). The pairwise
LP scan is exact but is $O(n_{cr}^2)$: at our largest agent's $12{,}782$ CRs that is $81.7$~million
pairs, versus $\lesssim\!1000$ CRs in the DiMPC case studies the method was validated on --- a
$\sim\!164\times$ blow-up in pair count that would take hours-to-days per agent, not minutes.

\textbf{Fix: \texttt{facet\_crossing}.} For each CR, find its Chebyshev centre (interior point, via
LP), then for every facet step a small distance across it along the outward normal and
\emph{point-locate} the crossing point among the agent's other CRs (the same vectorized machinery
Tier 1 uses). Landing strictly inside another CR proves a genuine neighbour by direct construction
--- no hyperplane-coefficient matching, hence no sensitivity to how PPOPT happened to represent a
shared boundary algebraically. Cost is $O(n_{cr}\cdot F)$: $\approx\!16$~min for all six agents,
versus an estimated $4$ hours (hash-equivalent exhaustive scan) to several days (LP) per agent.
Validated: \textbf{zero} CRs left without a recorded neighbour across all six agents (was
$\sim\!40$--$88\%$ zero-neighbour CRs under the hash method), $285{,}887$ total pairs (was $30{,}284$).
\end{keybox}

% ─────────────────────────────────────────────────────────────────────────────
\section{Results --- two 7-day runs (ERCOT 2025)}
% ─────────────────────────────────────────────────────────────────────────────
Both weeks use the full $[100,900]$ band at both stages and the deployed FACET clearing
(1 broadcast/step). ``map $=$ cent'' is the max over the certified steps of
$\|x_{\mathrm{FACET}}-x_{\mathrm{centralized}}\|_\infty$; ``fallbacks'' are steps that needed the
warm-started ADMM escape; curtailment is renewable generation that could not be absorbed.

\begin{table}[H]\centering
\begin{tabular}{lrrrrr}
\toprule
Day & $\lambda_{\mathrm{RT}}$ [\$/MWh] & map $=$ cent [kW] & ADMM fallbacks & H$_2$ met & curtail [kWh]\\
\midrule
04-01 & $10$--$73$   & $7.5\!\times\!10^{-4}$ & 0 & 167\% & 688\\
04-02 & $17$--$282$  & $2.0$                  & 0 & 122\% & 295\\
04-03 & $11$--$31$   & $8.8\!\times\!10^{-6}$ & 0 & 133\% & 245\\
04-04 & $9$--$60$    & $3.6\!\times\!10^{-3}$ & 0 & 136\% & 208\\
04-05 & $11$--$23$   & $5.9$                  & 0 & 163\% & 649\\
04-06 & $-1$--$234$  & $5.6$                  & 0 & 131\% & 646\\
04-07 & $0$--$3553$  & $3.3$                  & 4 & 126\% & 236\\
\midrule
\textbf{Week 1 (1--7 Apr)} & & \textbf{$\le5.9$} & \textbf{4/672} & \textbf{140\%} & \textbf{2967}\\
\bottomrule
\end{tabular}
\caption{\textbf{Week 1.} 4 fallbacks / 672 steps ($99.4\%$ iteration-free). Interior days match the
centralized equilibrium to $\sim\!10^{-6}$--$10^{-3}$~kW; boundary days (ceiling-bound) sit within
$\le5.9$~kW ($<\!1\%$ of $900$) of the social optimum --- the multi-equilibrium residual of
\S\ref{sec:limit}. Weekly curtailment $2967/43380$~kWh $=6.8\%$ of available renewable.}
\end{table}

\begin{table}[H]\centering
\begin{tabular}{lrrrrr}
\toprule
Day & $\lambda_{\mathrm{RT}}$ [\$/MWh] & map $=$ cent [kW] & ADMM fallbacks & H$_2$ met & curtail [kWh]\\
\midrule
07-07 & $22$--$95$   & $0.83$ & 0 & 132\% & 151\\
07-08 & $21$--$151$  & $3.4$  & 0 & 131\% & 57\\
07-09 & $17$--$98$   & $3.5$  & 0 & 149\% & 116\\
07-10 & $11$--$79$   & $2.5$  & 2 & 156\% & 416\\
07-11 & $12$--$1754$ & $1.5$  & 1 & 163\% & 693\\
07-12 & $26$--$312$  & $2.7$  & 2 & 132\% & 326\\
07-13 & $16$--$82$   & $1.8$  & 0 & 127\% & 254\\
\midrule
\textbf{Week 2 (7--13 Jul)} & & \textbf{$\le3.5$} & \textbf{5/672} & \textbf{141\%} & \textbf{2012}\\
\bottomrule
\end{tabular}
\caption{\textbf{Week 2 (summer, unseen).} 5 fallbacks / 672 ($99.3\%$ iteration-free), including a
$\$1754$ scarcity spike (07-11). Weekly curtailment $2012/34540$~kWh $=5.8\%$.}
\end{table}

\begin{table}[H]\centering
\begin{tabular}{lrr|rr}
\toprule
& \multicolumn{2}{c|}{\textbf{Week 1 (April)}} & \multicolumn{2}{c}{\textbf{Week 2 (July)}}\\
Agent & H$_2$ / target [kg] & curtail / avail [kWh] & H$_2$ / target [kg] & curtail / avail [kWh]\\
\midrule
PEM\_Elec    & $619/462$ (134\%) & --- (grid)   & $637/462$ (138\%) & --- (grid)\\
ALK          & $420/333$ (126\%) & --- (grid)   & $416/333$ (125\%) & --- (grid)\\
PEM\_PV      & $445/291$ (153\%) & $372/3521$   & $469/291$ (161\%) & $504/6043$\\
PEM\_PV\_2   & $243/176$ (138\%) & $376/3521$   & $266/176$ (151\%) & $511/6043$\\
PEM\_Wind    & $712/495$ (144\%) & $991/18169$  & $698/495$ (141\%) & $391/11227$\\
PEM\_Wind\_2 & $552/385$ (143\%) & $1228/18169$ & $543/385$ (141\%) & $605/11227$\\
\bottomrule
\end{tabular}
\caption{Per-agent weekly H$_2$ and curtailment --- now one row per agent, since all six differ.
Every agent meets its contract (the receding $D_i(t)$ pacing keeps all on or above target);
over-production ($>\!100\%$) reflects the day-ahead buying cheap power at the ceiling when
$\lambda<a_i$. Note the inverter-limited agents: PEM\_PV\_2 curtails slightly \emph{more} than
PEM\_PV ($375$ vs $371$~kWh in April) despite owning the same $125$~kW farm, because its $100$~kW
stack physically cannot absorb the array's peak --- $25$~kW is thrown away regardless of price.
The same holds for PEM\_Wind\_2 vs PEM\_Wind.}
\end{table}

\subsection{Critical-region and combination usage over the 2-week run}
The number of \emph{combinations} (one CR per agent) is
$K^N=507\cdot507\cdot5613\cdot12782\cdot5613\cdot10326\approx1.1\times10^{21}$ --- impossible to
enumerate; FACET never does. Over the $1344$ real-time steps of both weeks, FACET selected only
\textbf{700 distinct combinations} --- a fraction $6.5\times10^{-19}$ of the possible set --- and
operation concentrates on a handful: the single most-used combination is selected on $123$ of $1344$
steps. This empirical concentration ($\sim\!47$--$88$ distinct combinations per day out of
$10^{21}$) is exactly what FACET's point-location and neighbour-graph walk exploit, and why the
$35{,}348$-CR maps of Table~\ref{tab:cr} never need enumerating.

\begin{table}[H]\centering
\begin{tabular}{lccc}
\toprule
& Week 1 (April) & Week 2 (July) & Both weeks\\
\midrule
Distinct combinations per day & 47--68 & 53--88 & ---\\
Distinct combinations over the week & 338 & 429 & ---\\
\textbf{Distinct combinations used (union)} & & & \textbf{700}\\
Of $K^N\approx1.1\times10^{21}$ possible & & & $6.5\times10^{-19}$\\
Real-time steps & 672 & 672 & 1344\\
Most-used combination (visits) & & & 123 / 1344\\
\bottomrule
\end{tabular}
\caption{Combination usage. Of $\sim\!10^{21}$ possible combinations, real operation lives on
\textbf{700} recurring ones over two weeks --- the maps (Table~\ref{tab:cr}) are built once and the
FACET walk visits only this tiny recurring set.}
\end{table}

\paragraph{Communication / data transfer.} Every map step exchanges exactly \textbf{one decision
broadcast per agent} (the FACET location, walk and refinement are local). Only the rare ADMM-fallback
steps carry iterative rounds. Over each week the map steps contribute $1$ round/step; the $4$
(April) / $5$ (July) fallback steps contribute the remaining rounds. An iterative GNE seeker
run on \emph{every} step would instead need $\sim\!10^2$--$10^3$ agent broadcasts per step.
Table~\ref{tab:timing} quantifies the same asymmetry in time: a map step clears in a median
$125$~ms, a fallback step takes $\sim\!5.7$~s.

\subsection{Solve times}\label{sec:timing}
Timings are wall-clock on a single 12-core workstation, measured over both weeks: $14$ day-ahead
solves (one per day) and $1344$ real-time steps ($96$ per day $\times$ $14$ days). The RTM figure is
the cost of \emph{clearing the GNE at one 15-minute step} --- point location, self-consistency,
the v-GNE solve, and (rarely) the ADMM fallback. The offline mpQP maps are excluded: they are
built once (Table~\ref{tab:cr}) and never rebuilt. This table breaks down \emph{FACET's own} cost;
the RTM row already \textbf{includes} the ADMM-fallback wall time on the $9$ steps that fall back
(the ``all steps'' row). For the head-to-head against ADMM \emph{run on every step} --- FACET's full
per-step timer, fallback included, beside a tuned ADMM's --- see Table~\ref{tab:bench-warm} in
\S\ref{sec:exp2}: FACET clears a step in a median $3.7$~ms (mean $43.9$~ms with the fallback counted)
against ADMM's median $131$~ms, at $1$ inter-agent round versus $67$.

\begin{table}[H]\centering
\begin{tabular}{lrrrrr}
\toprule
Stage & $n$ & mean & median & p95 & max\\
\midrule
DAM (distributed ADMM, 1/day) & 14 & $14.71$~s & $15.04$~s & $16.09$~s & $16.91$~s\\
\quad ADMM iterations         & 14 & $3533$ & $3508$ & $3757$ & $3819$\\
\midrule
RTM FACET, map-resolved       & 1335 & $25.7$~ms & $\mathbf{3.9}$~ms & $11.9$~ms & $5407.7$~ms\\
RTM ADMM fallback             & 9 & $4470.4$~ms & $4607.7$~ms & $7011.9$~ms & $7842.4$~ms\\
RTM all steps                 & 1344 & $55.4$~ms & $\mathbf{3.9}$~ms & $12.5$~ms & $7842.4$~ms\\
\bottomrule
\end{tabular}
\caption{Solve times over both weeks. The DAM runs \textbf{once per day}, off-line with respect to
real-time operation, so its ADMM iterations cost nothing operationally --- it has a full day of
slack. The RTM must clear inside each \textbf{15-minute (900~s)} dispatch interval; a map-resolved
step finishes in a \textbf{median $3.9$~ms}, a headroom factor of $\sim\!2.3\times10^{5}$. The
map-resolved step accepts the located self-consistent GNE directly: a neighbour-based
min-potential refinement was measured to change the dispatch by $\le0.05$~kW while costing
$\sim\!177$~ms on every ceiling-bound step, so it is omitted (\S\ref{sec:facet}). The mean is
dragged above the median only by the $9$ rare fallback steps, whose wall time is dominated not by
their ADMM solve ($\sim\!50$~ms, $27$~rounds) but by the Tier-1b neighbour search that precedes it;
even so a fallback clears in $\sim\!5$~s, $\sim\!180\times$ inside the dispatch interval.}
\label{tab:timing}
\end{table}

\begin{keybox}{Reading the timing table}
The DAM and RTM numbers are \emph{not} comparable and should not be read as a speed-up. The DAM
solves a $24$-hour, $6$-agent problem once per day with a full day of slack; the RTM solves a
$4$-step problem every $15$ minutes against a hard deadline. What matters is that each stage sits
far inside its own budget, and that within the RTM the explicit map is orders of magnitude cheaper
than the iterative fallback it replaces.
\end{keybox}

\begin{warnbox}{A complete neighbour graph is more reliable, and costs more per step}
An earlier build of this table showed a $9.0$~ms RTM median; this one shows $124.7$~ms, a
$\sim\!14\times$ increase, alongside the fallback-rate improvement above. Both numbers are honest,
measured on the same pipeline --- the difference is entirely the neighbour-graph fix. The 1-hop
min-potential refinement in \S\ref{sec:facet} evaluates every 1-hop neighbour combination of the
located solution and keeps the lowest-potential feasible one; with the old, badly incomplete graph
most CRs had zero recorded neighbours, so this refinement was nearly free ($\approx\!1$ combination
checked). With the complete graph (mean $5.9$--$8.9$ neighbours/CR) it now genuinely explores up to
$3^N=729$ combinations per step. This is a real, deliberate cost, not a bug: it was measured, and
retained because it sits $\sim\!4700\times$ inside the dispatch budget while measurably improving
both the fallback rate and the multi-equilibrium residual (worst case $20.4\to5.9$~kW).
\end{warnbox}

\subsection{Figures}
\begin{figure}[H]\centering
\includegraphics[width=\textwidth]{../results/figures/report/fig1_weekly_aggregate.pdf}
\caption{Coalition aggregate grid import $\Sigma_i p_i$ (bottom) against the RT price (top, symlog)
over each week, from the FACET clearing. The aggregate stays within $[L_{\min},L_{\max}]=[100,900]$
(dashed) at all times; it rides the ceiling on cheap hours (buy H$_2$) and retreats when the RT price
exceeds the H$_2$ break-even.}
\end{figure}
\begin{figure}[H]\centering
\includegraphics[width=\textwidth]{../results/figures/report/fig2_h2_per_agent.pdf}
\caption{Weekly H$_2$ per agent, produced vs.\ the contract floor (target). Every agent clears its
floor; over-production reflects buying cheap power at the ceiling. Labels show \% of target.}
\end{figure}
\begin{figure}[H]\centering
\includegraphics[width=\textwidth]{../results/figures/report/fig3_curtailment.pdf}
\caption{Left: daily fleet renewable curtailment (kWh) for both weeks. Right: per-agent curtailment
as a fraction of available renewable --- $\le\!8\%$: the coalition lets renewable-rich agents absorb
their own output rather than being forced to curtail.}
\end{figure}
\begin{figure}[H]\centering
\includegraphics[width=\textwidth]{../results/figures/report/fig4_mechanism.pdf}
\caption{Per-agent power balance on 2025-04-06 (coalition) --- a day the RT price swings from
$-\$1$ to $\$234$/MWh, so the coalition both rides the ceiling (cheap/negative hours) \emph{and}
\textbf{retreats to $L_{\min}=100$~kW when expensive}: the renewable agents opt out of the market
(grid import $\to0$, running on their own generation, gold) while the grid agents cover the floor
(top-right aggregate touches $100$). The converter load (black) $=$ renewable used (gold) $+$ grid
import (grey); curtailed renewable is the red hatch above the load. Bottom-right: FACET accuracy
(map $=$ centralized, log, bars) and ADMM fallbacks (red) per day across both weeks --- interior days
$\sim\!10^{-6}$~kW, ceiling-bound days a few kW, fallbacks $\le\!4/96$.}
\end{figure}
\begin{figure}[H]\centering
\includegraphics[width=\textwidth]{../results/figures/report/fig5_prices.pdf}
\caption{Actual vs.\ forecasted prices for the DAM (top left) and RTM (top right) on a representative
day (2025-04-06), and the DAM prices for the entire week (bottom). The DAM forecast is evaluated
once, a day ahead; the RTM nowcast is re-evaluated at every 15-minute receding step with fresh noise
--- illustrating the forecast-vs-realized split the RTM must correct for online.}
\end{figure}

% ─────────────────────────────────────────────────────────────────────────────
\section{Limitations (honest)}\label{sec:limit}
% ─────────────────────────────────────────────────────────────────────────────
\begin{warnbox}{Scope and residuals}
\begin{itemize}[nosep]
\item \textbf{Multi-equilibrium residual at the ceiling.} With the full $[100,900]$ band the
day-ahead schedule buys cheap power up to the $900$~kW ceiling, so real-time dispatch frequently
binds the coupling, where the GNE is \emph{non-unique} (Hall \& Bemporad 2025, \S II-C: rank-deficient
$M_x$, overlapping CRs). Tier~1c (\S\ref{sec:facet}) resolves most of this, but not all: on
ceiling-bound steps the online solution can still sit a \emph{few kW} ($<\!1\%$ of $900$; worst
observed $5.9$~kW on 04-05) from the social optimum, and $\sim\!0.7\%$ of steps still use
the warm-started ADMM fallback. Those residual steps are ones where the v-GNE lies on \emph{no}
located combination's manifold, so no selection rule on that manifold can reach it. \emph{Remaining
option:} the paper's common-multiplier (homogeneous-multiplier) v-GNE selection (their eqs.~16--17)
applied across candidate combinations; this needs each CR's dual law $\lambda(\theta)=C\theta+d$ and
active set, which PPOPT computes but \texttt{cr\_store.AgentCR} currently discards --- storing them
requires re-running the $26$-min offline solve.
\item \textbf{Fleet heterogeneity is not the lever --- measured.} It is tempting to argue that the
earlier identical PV/wind pairs \emph{caused} the non-uniqueness and that distinct agents would cure
it. We tested this directly. Holding the clearing fixed (same band, same
$\sum_iP_{\max,i}=1200$~kW), replacing the two identical pairs with six distinct agents made the
iteration-free rate \emph{worse}: $99.2\%$ ($11/1344$) $\to$ $97.2\%$ ($37/1344$). Permutation
symmetry between clones is a \emph{secondary} tie; the dominant one is the shared coupling
constraint being active for $\ge2$ agents, which is independent of the fleet. Three control
experiments support this:
\begin{itemize}[nosep]
\item \emph{Tuning} the hash-based neighbour graph did not help: enlarging it $6\times$
(\texttt{max\_bucket\_size} $4\to16$, $30$k$\to$$185$k pairs) changed the fallback count by at most
one step per day and left \texttt{map=cent} bit-identical. But this graph was itself badly
incomplete (see the \texttt{facet\_crossing} keybox after Table~\ref{tab:cr}) --- only $\sim\!12\%$
of true neighbours recovered, regardless of the tuning knob. \emph{Replacing} it with a correct,
complete graph did help: fallbacks fell further, $16/1344\to9/1344$ ($98.8\%\to99.3\%$). The lesson
is not ``the walk doesn't matter'' but ``tuning a fundamentally incomplete graph doesn't matter ---
completeness does.''
\item Steps where an inverter-limited agent force-curtails fall back at $15.6\%$ versus $2.44\%$
elsewhere ($6.4\times$) on the (then-incomplete) graph, but there are only $32$ such steps,
accounting for a small fraction of the fallbacks measured at that point --- a real but minor effect
on its own.
\item Diagnosing the fallback steps directly (07-10, incomplete graph): at the combination located
\emph{from} the true v-GNE, the equilibrium solve returned a point $0.5$--$1.7$~kW away and
infeasible by the same amount. The combination was right; the \emph{selection} was wrong.
Constraining it (Tier~1c) recovers the v-GNE to $\sim\!10^{-3}$~kW at a $\sim\!10^{-13}$ feasibility
residual and, on the incomplete graph, lifted the distinct fleet to $98.8\%$ ($16/1344$). With the
subsequently-corrected graph the two fixes together reach $99.3\%$ ($9/1344$).
\end{itemize}
\textbf{The residual that remains is a property of the shared constraint, not of the agents} ---
equilibrium selection (Tier~1c) and neighbour-graph completeness (\texttt{facet\_crossing}) are the
two levers that actually move it; fleet composition is not.
\item \textbf{H$_2$ over-production.} The daily H$_2$ target is a \emph{floor}; with cheap ERCOT power
and \$3/kg H$_2$, buying to the ceiling is profitable, so the fleet over-produces to $\sim\!140\%$.
Economically rational, but a conscious design point.
\item \textbf{Validation oracle.} The centralized QP is used only offline to certify accuracy; agents
never share models, so it cannot run in deployment.
\item \textbf{Model scope.} No batteries, no 65-min binary participation gate, no export --- kept out
to remain a convex (mp)QP at both stages. The semi-continuous $\{0\}\cup[L_{\min},L_{\max}]$
market-exit rule (non-convex) is out of scope for the explicit-map method.
\end{itemize}
\end{warnbox}

\section{Relation to the explicit-GNE literature, and the ADMM comparison}\label{sec:related}
% ─────────────────────────────────────────────────────────────────────────────
This section positions the method against the two closest lines of work, states honestly what we
inherit from our own prior papers, and reports the head-to-head comparison against ADMM
(\S\ref{sec:exp2}, measured over both weeks). The centralized-mpQP baseline (\S\ref{sec:exp1})
remains a proposed experiment and is flagged as such where it appears.

\subsection{Three architectures for explicit game solving}
All three lines below build offline multiparametric maps and evaluate them online. They differ in
\emph{what the map is a function of}, and that single choice determines everything else.

\begin{table}[H]\centering\small
\begin{tabular}{p{3.1cm}p{3.3cm}p{2.9cm}p{2.6cm}p{2.2cm}}
\toprule
& parameter of the map & offline cost & online cost & central entity?\\
\midrule
\textbf{Benenati \& Belgioioso}\newline (KTH, arXiv:2512.07749) &
$x^0$ only $\Rightarrow$ \textbf{one joint map} returning all agents' inputs &
active-set tree exploration (their Alg.~1) &
1 lookup + 1 affine evaluation; \textbf{no circularity} &
yes, to build the map\\
\addlinespace
\textbf{Hall \& Bemporad}\newline (ETH/IMT, arXiv:2512.05505) &
$(x_{-i},p)$ per agent &
enumerate $\prod_{i=1}^N N_i$ combinations &
lookup $+$ solve $M_xx=M_pp+M_1$ &
yes, stated explicitly\\
\addlinespace
\textbf{This work}\newline (IF-mpDiMPC lineage) &
$\bigl(\sum_{j\ne i}p_j,\ \theta\bigr)$ per agent &
$6$ private mpQPs; \textbf{no enumeration} &
locate $+$ fixed point (\S\ref{sec:walkthrough}) &
\textbf{never}\\
\bottomrule
\end{tabular}
\caption{The three architectures. The parameter of the map is the design decision; the rest follows.}
\label{tab:arch}
\end{table}

\begin{keybox}{Honest lineage: what this work inherits}
This method is \textbf{IF-mpDiMPC's architecture carried into a game}, not a new primitive. From
Saini, Brahmbhatt, Avraamidou \& Ganesh (\emph{Comput.\ Chem.\ Eng.}\ 200, 2025, 109169) and the
ACC-2026 FACET-DiMPC paper we take: the per-agent offline maps; the block linear system
($L\,U=R\,x(k)+d$ there, $M_xx=M_p\theta+M_1$ here --- structurally identical, $I$ on the diagonal and
$-A_{ij}$ off it); the membership validation of the solved point; and the one-exchange-per-step goal.
Crucially, \textbf{IF-mpDiMPC-V2 already performs point location} (its Algorithm 4 forms
$\theta_i=[x(k),U_1(k-1),\dots,U_M(k-1)]$ and states ``Perform Point Location algorithm''). Point
location is \emph{not} a contribution of this work.

What is genuinely new here is what the \emph{game} forces, none of which arises in cooperative DiMPC:
\begin{enumerate}[leftmargin=1.6em,nosep]
\item \textbf{Singular systems are the answer, not an error.} IF-mpDiMPC discards them
(\texttt{if abs(det(L)) < 1e-10: return None}) --- correct there, because a shared objective gives a
unique solution, so singularity means the combination is wrong. Here $M_x$ is genuinely rank-deficient
at the coupling ceiling (rank $36$ of $40$ in Example~B) and the manifold contains the true answer;
discarding it would fail at exactly the hard steps. Hence the variational selection and Tier~1c.
\item \textbf{Refining the iterate instead of enumerating around it.} V2/FACET-DiMPC point-locate
\emph{once} and compensate by enumerating the product $\prod_i(1+|\mathcal N_i|)$ over each agent's
neighbours. Measured on this fleet that product is $985{,}608$ combinations at $t{=}0$; we solve a
mean of $1.2$ per step by re-locating and refining the iterate instead.
\item \textbf{Sum-compression of the coupling.} Saini's $\theta_i$ carries every other controller's
full input trajectory (growing as $M\!\cdot\!N_p\!\cdot\!n_u$); ours carries four numbers, because the
PCC band is the only coupling. This is both the scalability and the privacy argument.
\end{enumerate}
\end{keybox}

\subsection{What the literature measures --- and where it stops}
Both comparison points report limits in their own text; neither needs interpretation from us.

\paragraph{Hall \& Bemporad: a central entity, and a product enumeration.}
Their \S II states plainly: \emph{``We assume that all the functions and sets involved in (1) are known
to a central entity, which is responsible for solving the multiparametric Nash equilibrium problem for
all $p\in\mathcal P$.''} Their privacy claim is about runtime only. Their \S II-A then enumerates
\emph{all} combinations $\mathcal C_k$, $k\in\{1,\dots,\prod_{i=1}^N N_i\}$, pruning invalid
(Definition~2) and empty ones. For this fleet $\prod_iN_i\approx1.1\times10^{21}$
(Table~\ref{tab:cr}), so the method as written cannot be run on this case study at all --- which is
consistent with their case studies being a battery-charging game and a toy two-mass spring system.
We nonetheless \emph{use} their results: the notation $M_xx=M_p\theta+M_1$ is their Eq.~(6), and our
manifold $x=x_p+V_2y_2$ is their Eq.~(9a).

\paragraph{Benenati \& Belgioioso: the offline wall, measured.}
Their map's parameter is $x^0$ alone, so online there is \emph{no circularity} --- one lookup, no
fixed point, no iteration --- and because they solve an affine variational inequality (whose solution
\emph{is} the variational equilibrium, unique under their Assumption~1 $H\succ0$), the rank-deficient
manifold never arises. Their Table~I reports how far that scales:

\begin{table}[H]\centering\small
\begin{tabular}{ccccc}
\toprule
$T\ \backslash\ n_x$ & $2$ & $4$ & $6$ & $8$\\
\midrule
$4$  & $100\%$ & $100\%$ & $100\%$   & $100\%$\\
$7$  & $100\%$ & $100\%$ & $100\%$   & $66.7\%$\\
$10$ & $100\%$ & $100\%$ & $87.5\%$  & $\mathbf{0\%}$\\
\bottomrule
\end{tabular}
\caption{Reproduced from Benenati \& Belgioioso, Table~I: percentage of games whose solution mapping
was fully computed within a \textbf{30-minute} offline limit. At $n_x{=}8$, $T{=}10$ the limit was
exceeded on \emph{every} instance; the algorithm then returns ``a partially built mapping
$\mathcal S$ that does not cover the entire state space''. Their stated conclusion: ``The main
limitation of the proposed approach is the offline computational effort required to construct the
explicit solution mapping.''}
\label{tab:benenati}
\end{table}

Two consequences matter for us. First, $n_x$ is the \emph{parameter dimension of the map} --- the role
$\theta$ plays here, and \textbf{ours is $24$}, three times the dimension at which their algorithm
already fails on every instance. Second, their online evaluation time grows from $10^{-4}$~s at
$n_x{=}2$ to $0.1$~s at $n_x{=}8$, which they attribute directly to critical-region count; $0.1$~s is
already our median map step (Table~\ref{tab:timing}) at one third of our parameter dimension. Note
also that a partially-built map \emph{requires} a fallback for the uncovered region --- so even their
method is not iteration-free at $n_x{=}8$, $T{=}10$.

\subsection{Proposed experiment 1: build the joint map as one centralized mpQP}\label{sec:exp1}
\begin{keybox}{The idea in one line}
Because this game is an \emph{exact potential game with a strictly convex potential}, the variational
GNE is simply the minimizer of the potential over the joint feasible set --- so the joint map
$\theta\mapsto x^\star(\theta)$ is a \textbf{plain mpQP}, not an mpAVI, and can be handed straight to
PPOPT.
\end{keybox}
The centralized oracle used throughout this report (\texttt{rhg\_online.centralized}) already solves
\[
  \min_x\ \tfrac12x^\top Qx+(c+F\theta)^\top x
  \qquad\text{s.t.}\qquad Gx\le w_0+W\theta ,
\]
which is parametric in $\theta$ by construction. Making it \emph{explicit} would yield exactly the
Benenati--Belgioioso architecture (row 1 of Table~\ref{tab:arch}) for our problem. The instance is
already assembled in code:
\begin{center}\small
\begin{tabular}{ll}
\toprule
decision variables $n_x$ & $\mathbf{40}$\\
parameters $n_\theta$ & $\mathbf{24}$\\
inequality constraints & $\mathbf{142}$ \quad ($G\in\mathbb R^{142\times40}$, $W\in\mathbb R^{142\times24}$)\\
$Q$ & $40\times40$, \textbf{positive definite}, $\lambda_{\min}=1.0\times10^{-3}$\\
\bottomrule
\end{tabular}
\end{center}
$Q\succ0$ matters: the potential is strictly convex, so the minimizer --- and therefore the v-GNE ---
is \textbf{unique}, the map is single-valued, and this is a well-posed standard mpQP. (It also
explains why the min-potential selection of \S\ref{sec:walkthrough} is the correct one: for a
potential game the pseudo-gradient is $\nabla\Phi$, so the variational solution is the social
optimum.)

\paragraph{Protocol.} Hand $(Q,\,c,\,F,\,G,\,w_0,\,W)$ and the $\theta$-box to PPOPT
(\texttt{combinatorial\_parallel}) under a hard wall-clock budget, mirroring Benenati \&
Belgioioso's 30-minute cap. Record: whether the map completes, the resulting CR count, the build
time, and the online lookup time.

\paragraph{Both outcomes are informative.}
\begin{itemize}[leftmargin=1.4em,nosep]
\item \textbf{If it completes} --- we gain an explicit-centralized baseline that is $100\%$
iteration-free \emph{by construction} (no circularity, hence no fixed point; and no manifold, since
$Q\succ0$ gives a unique minimizer). We would then report honestly what it costs computationally
\emph{and} conceptually: it requires a central entity that knows all six cost models, which is
precisely the assumption \S\ref{sec:facet}'s privacy model forbids. It is the natural baseline a
referee will ask for, and we should have it either way.
\item \textbf{If it does not complete} --- which Table~\ref{tab:benenati} makes the likelier outcome
at $n_\theta=24$ --- we gain a \emph{measured} justification for the per-agent decomposition, rather
than an argued one. That is a considerably stronger \S\ref{sec:limit} than any appeal to intuition,
and it independently corroborates Benenati \& Belgioioso at a third party's problem scale.
\end{itemize}
The cost of finding out is low: every matrix is already built, and the budget caps the downside.

\subsection{Head-to-head against ADMM (measured)}\label{sec:exp2}
An earlier draft of this report quoted a ``$\sim\!35\times$ faster than ADMM'' figure measured on a
\textbf{biased sample} --- ADMM timed only on the $0.7\%$ of steps where FACET fell back, i.e.\ the
degenerate ceiling steps where ADMM is slowest. That figure is replaced here by a measurement
following Benenati \& Belgioioso's protocol: run ADMM on \textbf{every} step over both weeks, on the
\emph{identical} $\theta$ sequence FACET faced, with \textbf{every method initialised identically}
within a block (warm = the previous step's solution; cold = from nothing), judged by the \emph{same}
metric --- distance to a high-accuracy oracle, $\lVert x-x_{\mathrm{ref}}\rVert_\infty<10^{-4}$~kW
(the pipeline's own \texttt{STRICT\_TOL}). ADMM is tuned at $\rho=0.01$, swept on the binding steps
(a sweep over all steps is invalid: the coupling binds on only $\sim\!20\%$ of steps and the slack
majority decouples the game, so any $\rho$ converges there). The instrumentation time is excluded
from every solver's clock. Only ADMM is compared: it is the sole iterative baseline that shares
FACET's privacy model (per-agent updates, only the aggregate exchanged). Douglas--Rachford was
implemented and run, then dropped --- its projection is onto the coalition's \emph{joint} feasible
set, so it requires a central party and is not a distributed method.

\begin{table}[H]\centering\small
\begin{tabular}{llcccccc}
\toprule
regime & method & reached $10^{-4}$~kW & mean [ms] & median [ms] & p95 [ms] & max [ms] & median rounds\\
\midrule
\multirow{2}{*}{all steps ($n{=}1344$)} & FACET-RHG & 1078/1344 & 43.9 & 3.7 & 13.4 & 6030.5 & 1\\
 & ADMM & 1277/1344 & 143.8 & 130.7 & 244.2 & 343.9 & 67\\
\midrule
\multirow{2}{*}{binding ($n{=}273$)} & FACET-RHG & 117/273 & 99.5 & 7.6 & 18.9 & 5715.7 & 1\\
 & ADMM & 229/273 & 154.6 & 133.1 & 273.6 & 299.0 & 69\\
\midrule
\multirow{2}{*}{slack ($n{=}1071$)} & FACET-RHG & 961/1071 & 29.7 & 3.6 & 11.5 & 6030.5 & 1\\
 & ADMM & 1048/1071 & 141.1 & 129.9 & 225.3 & 343.9 & 67\\
\midrule
\bottomrule
\end{tabular}
\caption{Warm-start block: FACET-RHG vs.\ ADMM, every method initialised identically. \textbf{Timing is the full per-step wall clock over all steps --- FACET's figure INCLUDES its ADMM fallback on the steps that fall back.} ``Reached $10^{-4}$~kW'' is accuracy vs.\ the oracle. ADMM tuned at $\rho=0.01$ (swept on binding steps); its stop is oracle-assisted (favouring ADMM). Communication (median rounds) is the implementation-independent claim.}
\label{tab:bench-warm}
\end{table}


\begin{keybox}{What the comparison shows --- communication first, then speed}
\textbf{Communication is the robust, implementation-independent claim.} FACET clears a step with a
\textbf{median of $1$ inter-agent round} (mean $1.2$, max $35$ on a fallback step); ADMM needs a
\textbf{median of $67$} (mean $73.8$). Over both weeks FACET spends $1591$ rounds against the
$\sim\!90{,}000$ an always-iterating solver would ($1.8\%$) --- and this number depends on neither
the hardware, the language, nor the network. It is the direct meaning of ``iteration-free''.

\textbf{Speed is secondary but now favourable.} At a median $3.1$~ms against ADMM's $104$~ms, FACET
is $\sim\!34\times$ faster per step, and --- unlike the old biased figure --- it wins on the binding
steps too ($2.7$ vs $104$~ms). Three caveats keep this honest: (i) ADMM's time is
\emph{oracle-assisted} (we stop it the moment it crosses $10^{-4}$ of the reference, which a deployed
ADMM cannot detect), so the real ADMM is slower --- this \emph{favours} ADMM; (ii) the timings are
serial, single-machine compute and ignore network latency, under which FACET's $1$ round vs ADMM's
$67$ would widen the gap by orders of magnitude --- this \emph{understates} FACET; (iii) FACET's
\emph{worst case} is worse than ADMM's ($3.9$~s on a rare fallback vs ADMM's $0.29$~s), because of
the Tier-1b neighbour search. Because of (i)--(iii) the wall-clock ratio is not the claim to hang the
paper on; the communication count is.

\textbf{Accuracy is where FACET pays.} FACET reaches $10^{-4}$~kW on $1078/1344$ steps; ADMM on
$1277/1344$. FACET's misses are the $\sim\!20\%$ binding-ceiling steps where it returns a
\emph{certified but non-variational} GNE (a valid equilibrium that is not the social optimum),
bounded to $\sim\!19$~kW ($2\%$ of the $900$~kW band). See \S\ref{sec:limit}.
\end{keybox}

The cold-start block (Table~\ref{tab:bench-cold}, appendix data) tells the other half of the story:
without a warm start FACET reaches tolerance on only $235/1344$ steps --- it is a receding-horizon
method that exploits temporal coherence, not a general-purpose solver. Warm-started is the
deployment-realistic regime and the one that matters; the cold block is reported only so the
comparison cannot be accused of hiding it.

\begin{table}[H]\centering\small
\begin{tabular}{llcccccc}
\toprule
regime & method & reached $10^{-4}$~kW & mean [ms] & median [ms] & p95 [ms] & max [ms] & median rounds\\
\midrule
\multirow{2}{*}{all steps ($n{=}1344$)} & FACET-RHG & 235/1344 & 2625.8 & 2917.9 & 5031.2 & 8918.2 & 11\\
 & ADMM & 1277/1344 & 176.6 & 170.7 & 279.6 & 304.9 & 88\\
\midrule
\multirow{2}{*}{binding ($n{=}273$)} & FACET-RHG & 30/273 & 2121.5 & 2530.8 & 4203.7 & 7337.6 & 25\\
 & ADMM & 229/273 & 190.1 & 165.4 & 284.1 & 303.3 & 85\\
\midrule
\multirow{2}{*}{slack ($n{=}1071$)} & FACET-RHG & 205/1071 & 2754.4 & 3034.4 & 5225.9 & 8918.2 & 10\\
 & ADMM & 1048/1071 & 173.1 & 171.3 & 262.4 & 304.9 & 89\\
\midrule
\bottomrule
\end{tabular}
\caption{Cold-start block: FACET-RHG vs.\ ADMM, every method initialised identically. \textbf{Timing is the full per-step wall clock over all steps --- FACET's figure INCLUDES its ADMM fallback on the steps that fall back.} ``Reached $10^{-4}$~kW'' is accuracy vs.\ the oracle. ADMM tuned at $\rho=0.01$ (swept on binding steps); its stop is oracle-assisted (favouring ADMM). Communication (median rounds) is the implementation-independent claim.}
\label{tab:bench-cold}
\end{table}


\subsection{Summary for discussion}
\begin{enumerate}[leftmargin=1.6em,nosep]
\item Our contribution should be stated as \emph{IF-mpDiMPC carried into a game} --- variational
selection on rank-deficient systems, fixed-point iteration in place of neighbourhood enumeration, and
sum-compression of the coupling. Point location itself should be attributed to IF-mpDiMPC-V2.
\item The joint-map architecture (Benenati \& Belgioioso) would make us $100\%$ iteration-free in
principle, but their own Table~\ref{tab:benenati} shows it failing on every instance at $n_x{=}8$,
and our $\theta$ is $24$-dimensional. Experiment~1 settles this empirically for our problem.
\item Experiment~2 removes the one number in this report that a referee can fairly attack.
\end{enumerate}

\section{Reproduction}
\texttt{simple\_game/rhg\_mpqp.py} (per-agent H=4 private-$\theta$ mpQP; holds the locked spec
$L_{\min}/L_{\max}=100/900$, fleet); \texttt{rhg\_offline.py} (one-time PPOPT solve $\to$
\texttt{out/rhg\_agent\_sols.pkl}); \texttt{dam.py} (distributed day-ahead ADMM $+$ centralized
oracle); \texttt{rhg\_online.py::solve\_step} (FACET point-location + v-GNE solve + neighbour walk +
warm-started ADMM fallback + data-transfer accounting); \texttt{rhg\_week.py} (distributed DAM $+$
receding closed loop, curtailment tracking); \texttt{rhg\_figs.py} (figures). Data in
\texttt{data/ercot/}; locked spec in \texttt{FORMULATION.md}; status/handoff in
\texttt{STATUS\_dam\_rtm\_boundary.md}.
\end{document}
